% Môn: Vật lý
% Lớp: 11
% Chương: Chương II. Sóng
% Bài: L11C2 Bài 10. Thực hành: Đo tần số của sóng âm
% Số câu: 162
% App đọc file này trực tiếp — sửa rồi Commit trên GitHub, bấm Đồng bộ GitHub .tex

% L11C2 Bài 10. Thực hành: Đo tần số của sóng âm

% ===== Câu 1 =====
% ID: L11C2B10-01-DS
% Mức: TH
\dangbt{Nhận biết dụng cụ và bố trí thí nghiệm đo tần số âm}
\begin{ex}
% ID: L11C2B10-01-DS
% Mức: TH
Chuẩn bị thí nghiệm đo tần số âm.
\choiceTF
{\True Cần nguồn âm ổn định.}
{\True Microphone đặt đủ gần để thu tín hiệu rõ.}
{\True Thiết bị hiển thị cần có trục thời gian.}
{Thước đo chiều dài là dụng cụ duy nhất bắt buộc.}
\loigiai{
\begin{itemchoice}
\item (Đúng) Cần nguồn âm ổn định. (Vì): Cần nguồn âm ổn định. b) Đ: Microphone đặt đủ gần để thu tín hiệu rõ. c) Đ: Thiết bị hiển thị cần có trục thời gian. d) S: Thước đo chiều dài là dụng cụ duy nhất bắt buộc
\item (Đúng) Microphone đặt đủ gần để thu tín hiệu rõ.
\item (Đúng) Thiết bị hiển thị cần có trục thời gian.
\item (Sai) Thước đo chiều dài là dụng cụ duy nhất bắt buộc.
\end{itemchoice}
}
\end{ex}

% ===== Câu 2 =====
% ID: L11C2B10-01-TL
% Mức: TH
\begin{ex}
% ID: L11C2B10-01-TL
% Mức: TH
Trình bày cách bố trí thí nghiệm đo tần số của một âm thoa bằng microphone và dao động kí.
\choiceTF

\loigiai{
Đặt microphone gần âm thoa, nối với dao động kí/phần mềm, kích thích âm thoa, điều chỉnh thang thời gian để tín hiệu rõ rồi đo chu kì và tính $f=1/T$.
}
\end{ex}

% ===== Câu 3 =====
% ID: L11C2B10-01-TLN
% Mức: TH
\begin{ex}
% ID: L11C2B10-01-TLN
% Mức: TH
Trên màn hình, 5 chu kì chiếm 10 ms. Chu kì tính theo ms bằng bao nhiêu?
\shortans{2}
\loigiai{
Xác định đúng dữ kiện và đơn vị, áp dụng hệ thức của dạng bài rồi tính được kết quả 2.
}
\end{ex}

% ===== Câu 4 =====
% ID: L11C2B10-01-TN
% Mức: NB
\begin{ex}
% ID: L11C2B10-01-TN
% Mức: NB
Thiết bị dùng để thu âm và biến dao động âm thành tín hiệu điện là
\choice
{\True microphone}
{thước mét}
{nhiệt kế}
{lực kế}
\loigiai{
Đáp án A. microphone. Đối chiếu định nghĩa, hệ thức hoặc dữ kiện của bài toán cho kết luận trên.
}
\end{ex}

% ===== Câu 5 =====
% ID: L11C2B10-02-DS
% Mức: TH
\begin{ex}
% ID: L11C2B10-02-DS
% Mức: TH
Khi bố trí microphone và nguồn âm.
\choiceTF
{\True Tránh tiếng ồn nền giúp tín hiệu rõ hơn.}
{\True Không để tín hiệu vượt thang đo.}
{\True Nên giữ nguồn âm ổn định trong mỗi lần đo.}
{Di chuyển microphone tùy ý trong lúc ghi luôn làm kết quả chính xác hơn.}
\loigiai{
\begin{itemchoice}
\item (Đúng) Tránh tiếng ồn nền giúp tín hiệu rõ hơn. (Vì): Tránh tiếng ồn nền giúp tín hiệu rõ hơn. b) Đ: Không để tín hiệu vượt thang đo. c) Đ: Nên giữ nguồn âm ổn định trong mỗi lần đo. d) S: Di chuyển microphone tùy ý trong lúc ghi luôn làm kết quả chính xác hơn
\item (Đúng) Không để tín hiệu vượt thang đo.
\item (Đúng) Nên giữ nguồn âm ổn định trong mỗi lần đo.
\item (Sai) Di chuyển microphone tùy ý trong lúc ghi luôn làm kết quả chính xác hơn.
\end{itemchoice}
}
\end{ex}

% ===== Câu 6 =====
% ID: L11C2B10-02-TL
% Mức: TH
\begin{ex}
% ID: L11C2B10-02-TL
% Mức: TH
Giải thích vì sao nên đo thời gian của nhiều chu kì thay vì chỉ một chu kì.
\choiceTF

\loigiai{
Cùng một độ không đảm bảo khi đọc thời gian, đo khoảng dài hơn làm sai số tương đối nhỏ hơn; sau đó chia cho số chu kì.
}
\end{ex}

% ===== Câu 7 =====
% ID: L11C2B10-02-TLN
% Mức: TH
\begin{ex}
% ID: L11C2B10-02-TLN
% Mức: TH
Tín hiệu có chu kì 4 ms. Tần số tính theo Hz bằng bao nhiêu?
\shortans{250}
\loigiai{
Xác định đúng dữ kiện và đơn vị, áp dụng hệ thức của dạng bài rồi tính được kết quả 250.
}
\end{ex}

% ===== Câu 8 =====
% ID: L11C2B10-02-TN
% Mức: NB
\begin{ex}
% ID: L11C2B10-02-TN
% Mức: NB
Để quan sát dạng tín hiệu âm theo thời gian cần dùng
\choice
{la bàn}
{\True dao động kí hoặc phần mềm ghi âm}
{cân điện tử}
{ampe kế một chiều}
\loigiai{
Đáp án B. dao động kí hoặc phần mềm ghi âm. Đối chiếu định nghĩa, hệ thức hoặc dữ kiện của bài toán cho kết luận trên.
}
\end{ex}

% ===== Câu 9 =====
% ID: L11C2B10-03-DS
% Mức: TH
\begin{ex}
% ID: L11C2B10-03-DS
% Mức: TH
Quy trình đọc đồ thị âm.
\choiceTF
{\True Chọn hai đỉnh cùng chiều liên tiếp để đo một chu kì.}
{\True Có thể đo thời gian của nhiều chu kì rồi chia số chu kì.}
{\True Tần số bằng nghịch đảo chu kì.}
{Biên độ đồ thị quyết định trực tiếp tần số.}
\loigiai{
\begin{itemchoice}
\item (Đúng) Chọn hai đỉnh cùng chiều liên tiếp để đo một chu kì. (Vì): Chọn hai đỉnh cùng chiều liên tiếp để đo một chu kì. b) Đ: Có thể đo thời gian của nhiều chu kì rồi chia số chu kì. c) Đ: Tần số bằng nghịch đảo chu kì. d) S: Biên độ đồ thị quyết định trực tiếp tần số
\item (Đúng) Có thể đo thời gian của nhiều chu kì rồi chia số chu kì.
\item (Đúng) Tần số bằng nghịch đảo chu kì.
\item (Sai) Biên độ đồ thị quyết định trực tiếp tần số.
\end{itemchoice}
}
\end{ex}

% ===== Câu 10 =====
% ID: L11C2B10-03-TL
% Mức: VD
\begin{ex}
% ID: L11C2B10-03-TL
% Mức: VD
Nêu ba lưu ý an toàn và kĩ thuật khi thực hiện thí nghiệm âm.
\choiceTF

\loigiai{
Giữ âm lượng an toàn, nối thiết bị đúng cổng/thang đo, tránh va đập microphone và giảm tiếng ồn nền.
}
\end{ex}

% ===== Câu 11 =====
% ID: L11C2B10-03-TLN
% Mức: TH
\begin{ex}
% ID: L11C2B10-03-TLN
% Mức: TH
8 chu kì kéo dài 16 ms. Tần số tính theo Hz bằng bao nhiêu?
\shortans{500}
\loigiai{
Xác định đúng dữ kiện và đơn vị, áp dụng hệ thức của dạng bài rồi tính được kết quả 500.
}
\end{ex}

% ===== Câu 12 =====
% ID: L11C2B10-03-TN
% Mức: TH
\begin{ex}
% ID: L11C2B10-03-TN
% Mức: TH
Khi đo tần số âm, đại lượng cần đọc trực tiếp trên đồ thị là
\choice
{khối lượng nguồn}
{cường độ dòng điện}
{\True chu kì hoặc thời gian của nhiều chu kì}
{bước sóng trong chân không}
\loigiai{
Đáp án C. chu kì hoặc thời gian của nhiều chu kì. Đối chiếu định nghĩa, hệ thức hoặc dữ kiện của bài toán cho kết luận trên.
}
\end{ex}

% ===== Câu 13 =====
% ID: L11C2B10-04-DS
% Mức: VD
\begin{ex}
% ID: L11C2B10-04-DS
% Mức: VD
Xét các nhận định tổng hợp thuộc dạng “Nhận biết dụng cụ và bố trí thí nghiệm đo tần số âm”.
\choiceTF
{\True microphone.}
{\True dao động kí hoặc phần mềm ghi âm.}
{\True chu kì hoặc thời gian của nhiều chu kì.}
{thước mét.}
\loigiai{
\begin{itemchoice}
\item (Đúng) microphone. (Vì): microphone. b) Đ: dao động kí hoặc phần mềm ghi âm. c) Đ: chu kì hoặc thời gian của nhiều chu kì. d) S: thước mét
\item (Đúng) dao động kí hoặc phần mềm ghi âm.
\item (Đúng) chu kì hoặc thời gian của nhiều chu kì.
\item (Sai) thước mét.
\end{itemchoice}
}
\end{ex}

% ===== Câu 14 =====
% ID: L11C2B10-04-TL
% Mức: VD
\begin{ex}
% ID: L11C2B10-04-TL
% Mức: VD
Phân tích các bước lập luận và chỉ ra điều kiện áp dụng trong bài toán: Trình bày cách bố trí thí nghiệm đo tần số của một âm thoa bằng microphone và dao động kí.
\choiceTF

\loigiai{
Đặt microphone gần âm thoa, nối với dao động kí/phần mềm, kích thích âm thoa, điều chỉnh thang thời gian để tín hiệu rõ rồi đo chu kì và tính $f=1/T$. Cần nêu rõ điều kiện áp dụng, đổi đơn vị thống nhất và kiểm tra ý nghĩa vật lí của kết quả.
}
\end{ex}

% ===== Câu 15 =====
% ID: L11C2B10-04-TLN
% Mức: VD
\begin{ex}
% ID: L11C2B10-04-TLN
% Mức: VD
Từ kết quả của tình huống sau, nếu đại lượng cần tìm được nhân với hệ số 2 thì giá trị mới bằng bao nhiêu? Trên màn hình, 5 chu kì chiếm 10 ms. Chu kì tính theo ms bằng bao nhiêu?
\shortans{4}
\loigiai{
Xác định đúng dữ kiện và đơn vị, áp dụng hệ thức của dạng bài rồi tính được kết quả 4.
}
\end{ex}

% ===== Câu 16 =====
% ID: L11C2B10-04-TN
% Mức: TH
\begin{ex}
% ID: L11C2B10-04-TN
% Mức: TH
Trên màn hình, 5 chu kì chiếm 10 ms. Chu kì tính theo ms bằng bao nhiêu? Chọn giá trị đúng.
\choice
{4}
{1}
{3}
{\True 2}
\loigiai{
Đáp án D. 2. Đối chiếu định nghĩa, hệ thức hoặc dữ kiện của bài toán cho kết luận trên.
}
\end{ex}

% ===== Câu 17 =====
% ID: L11C2B10-05-DS
% Mức: VD
\begin{ex}
% ID: L11C2B10-05-DS
% Mức: VD
Đánh giá cách xử lí các dữ kiện định lượng của dạng “Nhận biết dụng cụ và bố trí thí nghiệm đo tần số âm”.
\choiceTF
{\True Kết quả của bài toán thứ nhất là 2.}
{\True Kết quả của bài toán thứ hai là 250.}
{\True Kết quả của bài toán thứ ba là 500.}
{Có thể bỏ qua hoàn toàn đơn vị và điều kiện vật lí khi tính toán.}
\loigiai{
\begin{itemchoice}
\item (Đúng) Kết quả của bài toán thứ nhất là 2. (Vì): Kết quả của bài toán thứ nhất là 2. b) Đ: Kết quả của bài toán thứ hai là 250. c) Đ: Kết quả của bài toán thứ ba là 500. d) S: Có thể bỏ qua hoàn toàn đơn vị và điều kiện vật lí khi tính toán
\item (Đúng) Kết quả của bài toán thứ hai là 250.
\item (Đúng) Kết quả của bài toán thứ ba là 500.
\item (Sai) Có thể bỏ qua hoàn toàn đơn vị và điều kiện vật lí khi tính toán.
\end{itemchoice}
}
\end{ex}

% ===== Câu 18 =====
% ID: L11C2B10-05-TL
% Mức: VD
\begin{ex}
% ID: L11C2B10-05-TL
% Mức: VD
Một học sinh giải bài sau nhưng bỏ qua điều kiện vật lí. Hãy sửa cách làm và trình bày lời giải đúng: Giải thích vì sao nên đo thời gian của nhiều chu kì thay vì chỉ một chu kì.
\choiceTF

\loigiai{
Cách giải đúng: Cùng một độ không đảm bảo khi đọc thời gian, đo khoảng dài hơn làm sai số tương đối nhỏ hơn; sau đó chia cho số chu kì. Sai lầm cần tránh là dùng hệ thức khi chưa xác định điều kiện biên, môi trường hoặc đơn vị.
}
\end{ex}

% ===== Câu 19 =====
% ID: L11C2B10-05-TLN
% Mức: VD
\begin{ex}
% ID: L11C2B10-05-TLN
% Mức: VD
Từ kết quả của tình huống sau, nếu đại lượng cần tìm được nhân với hệ số 0.5 thì giá trị mới bằng bao nhiêu? Tín hiệu có chu kì 4 ms. Tần số tính theo Hz bằng bao nhiêu?
\shortans{125}
\loigiai{
Xác định đúng dữ kiện và đơn vị, áp dụng hệ thức của dạng bài rồi tính được kết quả 125.
}
\end{ex}

% ===== Câu 20 =====
% ID: L11C2B10-05-TN
% Mức: TH
\begin{ex}
% ID: L11C2B10-05-TN
% Mức: TH
Tín hiệu có chu kì 4 ms. Tần số tính theo Hz bằng bao nhiêu? Chọn giá trị đúng.
\choice
{\True 250}
{500}
{125}
{251}
\loigiai{
Đáp án A. 250. Đối chiếu định nghĩa, hệ thức hoặc dữ kiện của bài toán cho kết luận trên.
}
\end{ex}

% ===== Câu 21 =====
% ID: L11C2B10-06-DS
% Mức: TH
\dangbt{Phân tích tín hiệu âm bằng dao động kí và phần mềm}
\begin{ex}
% ID: L11C2B10-06-DS
% Mức: TH
Trong chuyên đề “Phân tích tín hiệu âm bằng dao động kí và phần mềm”, Đồ thị cho 6 chu kì trong 12 ms.
\choiceTF
{\True Chu kì là 2 ms.}
{\True Tần số là $500\,\text{Hz}$.}
{\True Trong $1\,\text{s}$ có 500 dao động.}
{Tăng biên độ hiển thị làm tần số tăng.}
\loigiai{
\begin{itemchoice}
\item (Đúng) Chu kì là 2 ms. (Vì): Chu kì là 2 ms. b) Đ: Tần số là $500\,\text{Hz}$. c) Đ: Trong $1\,\text{s}$ có 500 dao động. d) S: Tăng biên độ hiển thị làm tần số tăng
\item (Đúng) Tần số là $500\,\text{Hz}$.
\item (Đúng) Trong $1\,\text{s}$ có 500 dao động.
\item (Sai) Tăng biên độ hiển thị làm tần số tăng.
\end{itemchoice}
}
\end{ex}

% ===== Câu 22 =====
% ID: L11C2B10-06-TL
% Mức: VD
\dangbt{Nhận biết dụng cụ và bố trí thí nghiệm đo tần số âm}
\begin{ex}
% ID: L11C2B10-06-TL
% Mức: VD
Mở rộng tình huống sau thành một ví dụ thực tế và trình bày phương pháp giải: Nêu ba lưu ý an toàn và kĩ thuật khi thực hiện thí nghiệm âm.
\choiceTF

\loigiai{
Giữ âm lượng an toàn, nối thiết bị đúng cổng/thang đo, tránh va đập microphone và giảm tiếng ồn nền. Có thể kiểm chứng bằng cách thay kết quả vào dữ kiện ban đầu và đối chiếu với hiện tượng thực tế.
}
\end{ex}

% ===== Câu 23 =====
% ID: L11C2B10-06-TLN
% Mức: VD
\begin{ex}
% ID: L11C2B10-06-TLN
% Mức: VD
Từ kết quả của tình huống sau, nếu đại lượng cần tìm được nhân với hệ số 1.5 thì giá trị mới bằng bao nhiêu? 8 chu kì kéo dài 16 ms. Tần số tính theo Hz bằng bao nhiêu?
\shortans{750}
\loigiai{
Xác định đúng dữ kiện và đơn vị, áp dụng hệ thức của dạng bài rồi tính được kết quả 750.
}
\end{ex}

% ===== Câu 24 =====
% ID: L11C2B10-06-TN
% Mức: TH
\begin{ex}
% ID: L11C2B10-06-TN
% Mức: TH
8 chu kì kéo dài 16 ms. Tần số tính theo Hz bằng bao nhiêu? Chọn giá trị đúng.
\choice
{501}
{\True 500}
{1000}
{250}
\loigiai{
Đáp án B. 500. Đối chiếu định nghĩa, hệ thức hoặc dữ kiện của bài toán cho kết luận trên.
}
\end{ex}

% ===== Câu 25 =====
% ID: L11C2B10-07-DS
% Mức: TH
\dangbt{Phân tích tín hiệu âm bằng dao động kí và phần mềm}
\begin{ex}
% ID: L11C2B10-07-DS
% Mức: TH
Trong chuyên đề “Phân tích tín hiệu âm bằng dao động kí và phần mềm”, Hai lần đo cho 5 chu kì lần lượt là 10,0 ms và 10,2 ms.
\choiceTF
{\True Có thể lấy trung bình thời gian trước khi tính chu kì.}
{\True Thời gian trung bình của 5 chu kì là 10,1 ms.}
{\True Chu kì trung bình là 2,02 ms.}
{Tần số trung bình chính xác bằng $202\,\text{Hz}$.}
\loigiai{
\begin{itemchoice}
\item (Đúng) Có thể lấy trung bình thời gian trước khi tính chu kì. (Vì): Có thể lấy trung bình thời gian trước khi tính chu kì. b) Đ: Thời gian trung bình của 5 chu kì là 10,1 ms. c) Đ: Chu kì trung bình là 2,02 ms. d) S: Tần số trung bình chính xác bằng $202\,\text{Hz}$
\item (Đúng) Thời gian trung bình của 5 chu kì là 10,1 ms.
\item (Đúng) Chu kì trung bình là 2,02 ms.
\item (Sai) Tần số trung bình chính xác bằng $202\,\text{Hz}$.
\end{itemchoice}
}
\end{ex}

% ===== Câu 26 =====
% ID: L11C2B10-07-TL
% Mức: TH
\begin{ex}
% ID: L11C2B10-07-TL
% Mức: TH
Trong chuyên đề “Phân tích tín hiệu âm bằng dao động kí và phần mềm”, Đồ thị âm cho thấy từ đỉnh thứ 2 đến đỉnh thứ 12 là 25 ms. Tính chu kì và tần số.
\choiceTF

\loigiai{
Giữa 11 đỉnh có 10 chu kì: $T=25/10=2,5$ ms; $f=1/(2,5\times10^{-3})=400$ Hz.
}
\end{ex}

% ===== Câu 27 =====
% ID: L11C2B10-07-TLN
% Mức: TH
\begin{ex}
% ID: L11C2B10-07-TLN
% Mức: TH
Trong chuyên đề “Phân tích tín hiệu âm bằng dao động kí và phần mềm”, 12 chu kì kéo dài 24 ms. Tần số tính theo Hz bằng bao nhiêu?
\shortans{500}
\loigiai{
Xác định đúng dữ kiện và đơn vị, áp dụng hệ thức của dạng bài rồi tính được kết quả 500.
}
\end{ex}

% ===== Câu 28 =====
% ID: L11C2B10-07-TN
% Mức: VD
\dangbt{Nhận biết dụng cụ và bố trí thí nghiệm đo tần số âm}
\begin{ex}
% ID: L11C2B10-07-TN
% Mức: VD
Kết luận nào phù hợp nhất khi giải bài toán sau: Trình bày cách bố trí thí nghiệm đo tần số của một âm thoa bằng microphone và dao động kí.
\choice
{Đại lượng cần tìm luôn bằng 0 trong mọi trường hợp.}
{Kết quả không phụ thuộc môi trường, tần số hoặc điều kiện biên.}
{\True Đặt microphone gần âm thoa, nối với dao động kí/phần mềm, kích thích âm thoa, điều chỉnh thang thời gian để tín hiệu rõ rồi đo chu kì và tính $f=1/T$.}
{Không cần sử dụng định nghĩa hay dữ kiện đã cho.}
\loigiai{
Đáp án C. Đặt microphone gần âm thoa, nối với dao động kí/phần mềm, kích thích âm thoa, điều chỉnh thang thời gian để tín hiệu rõ rồi đo chu kì và tính $f=1/T$.. Đối chiếu định nghĩa, hệ thức hoặc dữ kiện của bài toán cho kết luận trên.
}
\end{ex}

% ===== Câu 29 =====
% ID: L11C2B10-08-DS
% Mức: TH
\dangbt{Phân tích tín hiệu âm bằng dao động kí và phần mềm}
\begin{ex}
% ID: L11C2B10-08-DS
% Mức: TH
Trong chuyên đề “Phân tích tín hiệu âm bằng dao động kí và phần mềm”, Một tín hiệu âm có chu kì 1,25 ms.
\choiceTF
{\True Chu kì bằng $0,00125\,\text{s}$.}
{\True Tần số bằng $800\,\text{Hz}$.}
{\True Nếu màn hình hiển thị 4 chu kì thì khoảng thời gian là 5 ms.}
{Âm này là siêu âm.}
\loigiai{
\begin{itemchoice}
\item (Đúng) Chu kì bằng $0,00125\,\text{s}$. (Vì): Chu kì bằng $0,00125\,\text{s}$. b) Đ: Tần số bằng $800\,\text{Hz}$. c) Đ: Nếu màn hình hiển thị 4 chu kì thì khoảng thời gian là 5 ms. d) S: Âm này là siêu âm
\item (Đúng) Tần số bằng $800\,\text{Hz}$.
\item (Đúng) Nếu màn hình hiển thị 4 chu kì thì khoảng thời gian là 5 ms.
\item (Sai) Âm này là siêu âm.
\end{itemchoice}
}
\end{ex}

% ===== Câu 30 =====
% ID: L11C2B10-08-TL
% Mức: TH
\begin{ex}
% ID: L11C2B10-08-TL
% Mức: TH
Trong chuyên đề “Phân tích tín hiệu âm bằng dao động kí và phần mềm”, Ba lần đo thời gian của 10 chu kì là 19,8 ms; 20,0 ms; 20,2 ms. Tính tần số từ giá trị trung bình.
\choiceTF

\loigiai{
Thời gian trung bình của 10 chu kì là 20,0 ms; $T=2,00$ ms; $f=500$ Hz.
}
\end{ex}

% ===== Câu 31 =====
% ID: L11C2B10-08-TLN
% Mức: TH
\begin{ex}
% ID: L11C2B10-08-TLN
% Mức: TH
Trong chuyên đề “Phân tích tín hiệu âm bằng dao động kí và phần mềm”, Khoảng thời gian 4,5 ms chứa 3 chu kì. Tần số tính theo Hz, làm tròn đến hàng đơn vị, bằng bao nhiêu?
\shortans{667}
\loigiai{
Xác định đúng dữ kiện và đơn vị, áp dụng hệ thức của dạng bài rồi tính được kết quả 667.
}
\end{ex}

% ===== Câu 32 =====
% ID: L11C2B10-08-TN
% Mức: VD
\dangbt{Nhận biết dụng cụ và bố trí thí nghiệm đo tần số âm}
\begin{ex}
% ID: L11C2B10-08-TN
% Mức: VD
Kết luận nào phù hợp nhất khi giải bài toán sau: Giải thích vì sao nên đo thời gian của nhiều chu kì thay vì chỉ một chu kì.
\choice
{Không cần sử dụng định nghĩa hay dữ kiện đã cho.}
{Đại lượng cần tìm luôn bằng 0 trong mọi trường hợp.}
{Kết quả không phụ thuộc môi trường, tần số hoặc điều kiện biên.}
{\True Cùng một độ không đảm bảo khi đọc thời gian, đo khoảng dài hơn làm sai số tương đối nhỏ hơn}
\loigiai{
Đáp án D. Cùng một độ không đảm bảo khi đọc thời gian, đo khoảng dài hơn làm sai số tương đối nhỏ hơn. Đối chiếu định nghĩa, hệ thức hoặc dữ kiện của bài toán cho kết luận trên.
}
\end{ex}

% ===== Câu 33 =====
% ID: L11C2B10-09-DS
% Mức: VD
\dangbt{Phân tích tín hiệu âm bằng dao động kí và phần mềm}
\begin{ex}
% ID: L11C2B10-09-DS
% Mức: VD
Trong chuyên đề “Phân tích tín hiệu âm bằng dao động kí và phần mềm”, Xét các nhận định tổng hợp thuộc dạng “Phân tích tín hiệu âm bằng dao động kí và phần mềm”.
\choiceTF
{\True $500\,\text{Hz}$.}
{\True $400\,\text{Hz}$.}
{\True đo thời gian của nhiều chu kì.}
{$50\,\text{Hz}$.}
\loigiai{
\begin{itemchoice}
\item (Đúng) $500\,\text{Hz}$. (Vì): $500\,\text{Hz}$. b) Đ: $400\,\text{Hz}$. c) Đ: đo thời gian của nhiều chu kì. d) S: $50\,\text{Hz}$
\item (Đúng) $400\,\text{Hz}$.
\item (Đúng) đo thời gian của nhiều chu kì.
\item (Sai) $50\,\text{Hz}$.
\end{itemchoice}
}
\end{ex}

% ===== Câu 34 =====
% ID: L11C2B10-09-TL
% Mức: VD
\begin{ex}
% ID: L11C2B10-09-TL
% Mức: VD
Trong chuyên đề “Phân tích tín hiệu âm bằng dao động kí và phần mềm”, Mô tả cách nhận biết và loại bỏ một đoạn tín hiệu bị nhiễu trước khi tính tần số.
\choiceTF

\loigiai{
Chọn đoạn tuần hoàn ổn định, biên dạng lặp rõ; bỏ đoạn có xung bất thường hoặc mất chu kì, rồi đo nhiều chu kì trên phần còn lại.
}
\end{ex}

% ===== Câu 35 =====
% ID: L11C2B10-09-TLN
% Mức: TH
\begin{ex}
% ID: L11C2B10-09-TLN
% Mức: TH
Trong chuyên đề “Phân tích tín hiệu âm bằng dao động kí và phần mềm”, Hai đỉnh thứ nhất và thứ chín cách nhau 16 ms. Tần số tính theo Hz bằng bao nhiêu?
\shortans{500}
\loigiai{
Xác định đúng dữ kiện và đơn vị, áp dụng hệ thức của dạng bài rồi tính được kết quả 500.
}
\end{ex}

% ===== Câu 36 =====
% ID: L11C2B10-09-TN
% Mức: VD
\dangbt{Nhận biết dụng cụ và bố trí thí nghiệm đo tần số âm}
\begin{ex}
% ID: L11C2B10-09-TN
% Mức: VD
Kết luận nào phù hợp nhất khi giải bài toán sau: Nêu ba lưu ý an toàn và kĩ thuật khi thực hiện thí nghiệm âm.
\choice
{\True Giữ âm lượng an toàn, nối thiết bị đúng cổng/thang đo, tránh va đập microphone và giảm tiếng ồn nền.}
{Không cần sử dụng định nghĩa hay dữ kiện đã cho.}
{Đại lượng cần tìm luôn bằng 0 trong mọi trường hợp.}
{Kết quả không phụ thuộc môi trường, tần số hoặc điều kiện biên.}
\loigiai{
Đáp án A. Giữ âm lượng an toàn, nối thiết bị đúng cổng/thang đo, tránh va đập microphone và giảm tiếng ồn nền.. Đối chiếu định nghĩa, hệ thức hoặc dữ kiện của bài toán cho kết luận trên.
}
\end{ex}

% ===== Câu 37 =====
% ID: L11C2B10-10-DS
% Mức: VD
\dangbt{Phân tích tín hiệu âm bằng dao động kí và phần mềm}
\begin{ex}
% ID: L11C2B10-10-DS
% Mức: VD
Trong chuyên đề “Phân tích tín hiệu âm bằng dao động kí và phần mềm”, Đánh giá cách xử lí các dữ kiện định lượng của dạng “Phân tích tín hiệu âm bằng dao động kí và phần mềm”.
\choiceTF
{\True Kết quả của bài toán thứ nhất là 500.}
{\True Kết quả của bài toán thứ hai là 667.}
{\True Kết quả của bài toán thứ ba là 500.}
{Có thể bỏ qua hoàn toàn đơn vị và điều kiện vật lí khi tính toán.}
\loigiai{
\begin{itemchoice}
\item (Đúng) Kết quả của bài toán thứ nhất là 500. (Vì): Kết quả của bài toán thứ nhất là 500. b) Đ: Kết quả của bài toán thứ hai là 667. c) Đ: Kết quả của bài toán thứ ba là 500. d) S: Có thể bỏ qua hoàn toàn đơn vị và điều kiện vật lí khi tính toán
\item (Đúng) Kết quả của bài toán thứ hai là 667.
\item (Đúng) Kết quả của bài toán thứ ba là 500.
\item (Sai) Có thể bỏ qua hoàn toàn đơn vị và điều kiện vật lí khi tính toán.
\end{itemchoice}
}
\end{ex}

% ===== Câu 38 =====
% ID: L11C2B10-10-TL
% Mức: VD
\begin{ex}
% ID: L11C2B10-10-TL
% Mức: VD
Trong chuyên đề “Phân tích tín hiệu âm bằng dao động kí và phần mềm”, Phân tích các bước lập luận và chỉ ra điều kiện áp dụng trong bài toán: Đồ thị âm cho thấy từ đỉnh thứ 2 đến đỉnh thứ 12 là 25 ms. Tính chu kì và tần số.
\choiceTF

\loigiai{
Giữa 11 đỉnh có 10 chu kì: $T=25/10=2,5$ ms; $f=1/(2,5\times10^{-3})=400$ Hz. Cần nêu rõ điều kiện áp dụng, đổi đơn vị thống nhất và kiểm tra ý nghĩa vật lí của kết quả.
}
\end{ex}

% ===== Câu 39 =====
% ID: L11C2B10-10-TLN
% Mức: VD
\begin{ex}
% ID: L11C2B10-10-TLN
% Mức: VD
Trong chuyên đề “Phân tích tín hiệu âm bằng dao động kí và phần mềm”, Từ kết quả của tình huống sau, nếu đại lượng cần tìm được nhân với hệ số 2 thì giá trị mới bằng bao nhiêu? 12 chu kì kéo dài 24 ms. Tần số tính theo Hz bằng bao nhiêu?
\shortans{1000}
\loigiai{
Xác định đúng dữ kiện và đơn vị, áp dụng hệ thức của dạng bài rồi tính được kết quả 1000.
}
\end{ex}

% ===== Câu 40 =====
% ID: L11C2B10-10-TN
% Mức: VD
\dangbt{Nhận biết dụng cụ và bố trí thí nghiệm đo tần số âm}
\begin{ex}
% ID: L11C2B10-10-TN
% Mức: VD
Phương pháp phù hợp nhất cho dạng “Nhận biết dụng cụ và bố trí thí nghiệm đo tần số âm” là
\choice
{luôn coi mọi điểm dao động cùng pha}
{\True xác định dữ kiện, chọn đúng hệ thức hoặc định nghĩa, đổi đơn vị rồi kiểm tra kết quả}
{chỉ thay số mà không cần xét điều kiện áp dụng}
{bỏ qua đơn vị và chiều truyền sóng}
\loigiai{
Đáp án B. xác định dữ kiện, chọn đúng hệ thức hoặc định nghĩa, đổi đơn vị rồi kiểm tra kết quả. Đối chiếu định nghĩa, hệ thức hoặc dữ kiện của bài toán cho kết luận trên.
}
\end{ex}

% ===== Câu 41 =====
% ID: L11C2B10-11-DS
% Mức: TH
\dangbt{Thiết kế phương án thí nghiệm đo tần số của nguồn âm}
\begin{ex}
% ID: L11C2B10-11-DS
% Mức: TH
Trong chuyên đề “Thiết kế phương án thí nghiệm đo tần số của nguồn âm”, Chuẩn bị thí nghiệm đo tần số âm.
\choiceTF
{\True Cần nguồn âm ổn định.}
{\True Microphone đặt đủ gần để thu tín hiệu rõ.}
{\True Thiết bị hiển thị cần có trục thời gian.}
{Thước đo chiều dài là dụng cụ duy nhất bắt buộc.}
\loigiai{
\begin{itemchoice}
\item (Đúng) Cần nguồn âm ổn định. (Vì): Cần nguồn âm ổn định. b) Đ: Microphone đặt đủ gần để thu tín hiệu rõ. c) Đ: Thiết bị hiển thị cần có trục thời gian. d) S: Thước đo chiều dài là dụng cụ duy nhất bắt buộc
\item (Đúng) Microphone đặt đủ gần để thu tín hiệu rõ.
\item (Đúng) Thiết bị hiển thị cần có trục thời gian.
\item (Sai) Thước đo chiều dài là dụng cụ duy nhất bắt buộc.
\end{itemchoice}
}
\end{ex}

% ===== Câu 42 =====
% ID: L11C2B10-11-TL
% Mức: VD
\dangbt{Phân tích tín hiệu âm bằng dao động kí và phần mềm}
\begin{ex}
% ID: L11C2B10-11-TL
% Mức: VD
Trong chuyên đề “Phân tích tín hiệu âm bằng dao động kí và phần mềm”, Một học sinh giải bài sau nhưng bỏ qua điều kiện vật lí. Hãy sửa cách làm và trình bày lời giải đúng: Ba lần đo thời gian của 10 chu kì là 19,8 ms; 20,0 ms; 20,2 ms. Tính tần số từ giá trị trung bình.
\choiceTF

\loigiai{
Cách giải đúng: Thời gian trung bình của 10 chu kì là 20,0 ms; $T=2,00$ ms; $f=500$ Hz. Sai lầm cần tránh là dùng hệ thức khi chưa xác định điều kiện biên, môi trường hoặc đơn vị.
}
\end{ex}

% ===== Câu 43 =====
% ID: L11C2B10-11-TLN
% Mức: VD
\begin{ex}
% ID: L11C2B10-11-TLN
% Mức: VD
Trong chuyên đề “Phân tích tín hiệu âm bằng dao động kí và phần mềm”, Từ kết quả của tình huống sau, nếu đại lượng cần tìm được nhân với hệ số 0.5 thì giá trị mới bằng bao nhiêu? Khoảng thời gian 4,5 ms chứa 3 chu kì. Tần số tính theo Hz, làm tròn đến hàng đơn vị, bằng bao nhiêu?
\shortans{333.5}
\loigiai{
Xác định đúng dữ kiện và đơn vị, áp dụng hệ thức của dạng bài rồi tính được kết quả 333.5.
}
\end{ex}

% ===== Câu 44 =====
% ID: L11C2B10-11-TN
% Mức: NB
\begin{ex}
% ID: L11C2B10-11-TN
% Mức: NB
Trong chuyên đề “Phân tích tín hiệu âm bằng dao động kí và phần mềm”, Trên đồ thị, 10 chu kì kéo dài 20 ms. Tần số âm bằng
\choice
{\True $500\,\text{Hz}$}
{$50\,\text{Hz}$}
{$200\,\text{Hz}$}
{$2000\,\text{Hz}$}
\loigiai{
Đáp án A. $500\,\text{Hz}$. Đối chiếu định nghĩa, hệ thức hoặc dữ kiện của bài toán cho kết luận trên.
}
\end{ex}

% ===== Câu 45 =====
% ID: L11C2B10-12-DS
% Mức: TH
\dangbt{Thiết kế phương án thí nghiệm đo tần số của nguồn âm}
\begin{ex}
% ID: L11C2B10-12-DS
% Mức: TH
Trong chuyên đề “Thiết kế phương án thí nghiệm đo tần số của nguồn âm”, Khi bố trí microphone và nguồn âm.
\choiceTF
{\True Tránh tiếng ồn nền giúp tín hiệu rõ hơn.}
{\True Không để tín hiệu vượt thang đo.}
{\True Nên giữ nguồn âm ổn định trong mỗi lần đo.}
{Di chuyển microphone tùy ý trong lúc ghi luôn làm kết quả chính xác hơn.}
\loigiai{
\begin{itemchoice}
\item (Đúng) Tránh tiếng ồn nền giúp tín hiệu rõ hơn. (Vì): Tránh tiếng ồn nền giúp tín hiệu rõ hơn. b) Đ: Không để tín hiệu vượt thang đo. c) Đ: Nên giữ nguồn âm ổn định trong mỗi lần đo. d) S: Di chuyển microphone tùy ý trong lúc ghi luôn làm kết quả chính xác hơn
\item (Đúng) Không để tín hiệu vượt thang đo.
\item (Đúng) Nên giữ nguồn âm ổn định trong mỗi lần đo.
\item (Sai) Di chuyển microphone tùy ý trong lúc ghi luôn làm kết quả chính xác hơn.
\end{itemchoice}
}
\end{ex}

% ===== Câu 46 =====
% ID: L11C2B10-12-TL
% Mức: VD
\dangbt{Phân tích tín hiệu âm bằng dao động kí và phần mềm}
\begin{ex}
% ID: L11C2B10-12-TL
% Mức: VD
Trong chuyên đề “Phân tích tín hiệu âm bằng dao động kí và phần mềm”, Mở rộng tình huống sau thành một ví dụ thực tế và trình bày phương pháp giải: Mô tả cách nhận biết và loại bỏ một đoạn tín hiệu bị nhiễu trước khi tính tần số.
\choiceTF

\loigiai{
Chọn đoạn tuần hoàn ổn định, biên dạng lặp rõ; bỏ đoạn có xung bất thường hoặc mất chu kì, rồi đo nhiều chu kì trên phần còn lại. Có thể kiểm chứng bằng cách thay kết quả vào dữ kiện ban đầu và đối chiếu với hiện tượng thực tế.
}
\end{ex}

% ===== Câu 47 =====
% ID: L11C2B10-12-TLN
% Mức: VD
\begin{ex}
% ID: L11C2B10-12-TLN
% Mức: VD
Trong chuyên đề “Phân tích tín hiệu âm bằng dao động kí và phần mềm”, Từ kết quả của tình huống sau, nếu đại lượng cần tìm được nhân với hệ số 1.5 thì giá trị mới bằng bao nhiêu? Hai đỉnh thứ nhất và thứ chín cách nhau 16 ms. Tần số tính theo Hz bằng bao nhiêu?
\shortans{750}
\loigiai{
Xác định đúng dữ kiện và đơn vị, áp dụng hệ thức của dạng bài rồi tính được kết quả 750.
}
\end{ex}

% ===== Câu 48 =====
% ID: L11C2B10-12-TN
% Mức: NB
\begin{ex}
% ID: L11C2B10-12-TN
% Mức: NB
Trong chuyên đề “Phân tích tín hiệu âm bằng dao động kí và phần mềm”, Khoảng thời gian giữa hai đỉnh liên tiếp là 2,5 ms. Tần số bằng
\choice
{4 kHz}
{\True $400\,\text{Hz}$}
{$250\,\text{Hz}$}
{$40\,\text{Hz}$}
\loigiai{
Đáp án B. $400\,\text{Hz}$. Đối chiếu định nghĩa, hệ thức hoặc dữ kiện của bài toán cho kết luận trên.
}
\end{ex}

% ===== Câu 49 =====
% ID: L11C2B10-13-DS
% Mức: TH
\dangbt{Thiết kế phương án thí nghiệm đo tần số của nguồn âm}
\begin{ex}
% ID: L11C2B10-13-DS
% Mức: TH
Trong chuyên đề “Thiết kế phương án thí nghiệm đo tần số của nguồn âm”, Quy trình đọc đồ thị âm.
\choiceTF
{\True Chọn hai đỉnh cùng chiều liên tiếp để đo một chu kì.}
{\True Có thể đo thời gian của nhiều chu kì rồi chia số chu kì.}
{\True Tần số bằng nghịch đảo chu kì.}
{Biên độ đồ thị quyết định trực tiếp tần số.}
\loigiai{
\begin{itemchoice}
\item (Đúng) Chọn hai đỉnh cùng chiều liên tiếp để đo một chu kì. (Vì): Chọn hai đỉnh cùng chiều liên tiếp để đo một chu kì. b) Đ: Có thể đo thời gian của nhiều chu kì rồi chia số chu kì. c) Đ: Tần số bằng nghịch đảo chu kì. d) S: Biên độ đồ thị quyết định trực tiếp tần số
\item (Đúng) Có thể đo thời gian của nhiều chu kì rồi chia số chu kì.
\item (Đúng) Tần số bằng nghịch đảo chu kì.
\item (Sai) Biên độ đồ thị quyết định trực tiếp tần số.
\end{itemchoice}
}
\end{ex}

% ===== Câu 50 =====
% ID: L11C2B10-13-TL
% Mức: TH
\begin{ex}
% ID: L11C2B10-13-TL
% Mức: TH
Trong chuyên đề “Thiết kế phương án thí nghiệm đo tần số của nguồn âm”, Trình bày cách bố trí thí nghiệm đo tần số của một âm thoa bằng microphone và dao động kí.
\choiceTF

\loigiai{
Đặt microphone gần âm thoa, nối với dao động kí/phần mềm, kích thích âm thoa, điều chỉnh thang thời gian để tín hiệu rõ rồi đo chu kì và tính $f=1/T$.
}
\end{ex}

% ===== Câu 51 =====
% ID: L11C2B10-13-TLN
% Mức: TH
\begin{ex}
% ID: L11C2B10-13-TLN
% Mức: TH
Trong chuyên đề “Thiết kế phương án thí nghiệm đo tần số của nguồn âm”, Trên màn hình, 5 chu kì chiếm 10 ms. Chu kì tính theo ms bằng bao nhiêu?
\shortans{2}
\loigiai{
Xác định đúng dữ kiện và đơn vị, áp dụng hệ thức của dạng bài rồi tính được kết quả 2.
}
\end{ex}

% ===== Câu 52 =====
% ID: L11C2B10-13-TN
% Mức: TH
\dangbt{Phân tích tín hiệu âm bằng dao động kí và phần mềm}
\begin{ex}
% ID: L11C2B10-13-TN
% Mức: TH
Trong chuyên đề “Phân tích tín hiệu âm bằng dao động kí và phần mềm”, Để giảm sai số đọc chu kì từ đồ thị, nên
\choice
{đọc biên độ thay cho thời gian}
{thay đổi tần số khi đang đo}
{\True đo thời gian của nhiều chu kì}
{chỉ chọn một đoạn rất ngắn}
\loigiai{
Đáp án C. đo thời gian của nhiều chu kì. Đối chiếu định nghĩa, hệ thức hoặc dữ kiện của bài toán cho kết luận trên.
}
\end{ex}

% ===== Câu 53 =====
% ID: L11C2B10-14-DS
% Mức: VD
\dangbt{Thiết kế phương án thí nghiệm đo tần số của nguồn âm}
\begin{ex}
% ID: L11C2B10-14-DS
% Mức: VD
Trong chuyên đề “Thiết kế phương án thí nghiệm đo tần số của nguồn âm”, Xét các nhận định tổng hợp thuộc dạng “Thiết kế phương án thí nghiệm đo tần số của nguồn âm”.
\choiceTF
{\True microphone.}
{\True dao động kí hoặc phần mềm ghi âm.}
{\True chu kì hoặc thời gian của nhiều chu kì.}
{thước mét.}
\loigiai{
\begin{itemchoice}
\item (Đúng) microphone. (Vì): microphone. b) Đ: dao động kí hoặc phần mềm ghi âm. c) Đ: chu kì hoặc thời gian của nhiều chu kì. d) S: thước mét
\item (Đúng) dao động kí hoặc phần mềm ghi âm.
\item (Đúng) chu kì hoặc thời gian của nhiều chu kì.
\item (Sai) thước mét.
\end{itemchoice}
}
\end{ex}

% ===== Câu 54 =====
% ID: L11C2B10-14-TL
% Mức: TH
\begin{ex}
% ID: L11C2B10-14-TL
% Mức: TH
Trong chuyên đề “Thiết kế phương án thí nghiệm đo tần số của nguồn âm”, Giải thích vì sao nên đo thời gian của nhiều chu kì thay vì chỉ một chu kì.
\choiceTF

\loigiai{
Cùng một độ không đảm bảo khi đọc thời gian, đo khoảng dài hơn làm sai số tương đối nhỏ hơn; sau đó chia cho số chu kì.
}
\end{ex}

% ===== Câu 55 =====
% ID: L11C2B10-14-TLN
% Mức: TH
\begin{ex}
% ID: L11C2B10-14-TLN
% Mức: TH
Trong chuyên đề “Thiết kế phương án thí nghiệm đo tần số của nguồn âm”, Tín hiệu có chu kì 4 ms. Tần số tính theo Hz bằng bao nhiêu?
\shortans{250}
\loigiai{
Xác định đúng dữ kiện và đơn vị, áp dụng hệ thức của dạng bài rồi tính được kết quả 250.
}
\end{ex}

% ===== Câu 56 =====
% ID: L11C2B10-14-TN
% Mức: TH
\dangbt{Phân tích tín hiệu âm bằng dao động kí và phần mềm}
\begin{ex}
% ID: L11C2B10-14-TN
% Mức: TH
Trong chuyên đề “Phân tích tín hiệu âm bằng dao động kí và phần mềm”, 12 chu kì kéo dài 24 ms. Tần số tính theo Hz bằng bao nhiêu? Chọn giá trị đúng.
\choice
{1000}
{250}
{501}
{\True 500}
\loigiai{
Đáp án D. 500. Đối chiếu định nghĩa, hệ thức hoặc dữ kiện của bài toán cho kết luận trên.
}
\end{ex}

% ===== Câu 57 =====
% ID: L11C2B10-15-DS
% Mức: VD
\dangbt{Thiết kế phương án thí nghiệm đo tần số của nguồn âm}
\begin{ex}
% ID: L11C2B10-15-DS
% Mức: VD
Trong chuyên đề “Thiết kế phương án thí nghiệm đo tần số của nguồn âm”, Đánh giá cách xử lí các dữ kiện định lượng của dạng “Thiết kế phương án thí nghiệm đo tần số của nguồn âm”.
\choiceTF
{\True Kết quả của bài toán thứ nhất là 2.}
{\True Kết quả của bài toán thứ hai là 250.}
{\True Kết quả của bài toán thứ ba là 500.}
{Có thể bỏ qua hoàn toàn đơn vị và điều kiện vật lí khi tính toán.}
\loigiai{
\begin{itemchoice}
\item (Đúng) Kết quả của bài toán thứ nhất là 2. (Vì): Kết quả của bài toán thứ nhất là 2. b) Đ: Kết quả của bài toán thứ hai là 250. c) Đ: Kết quả của bài toán thứ ba là 500. d) S: Có thể bỏ qua hoàn toàn đơn vị và điều kiện vật lí khi tính toán
\item (Đúng) Kết quả của bài toán thứ hai là 250.
\item (Đúng) Kết quả của bài toán thứ ba là 500.
\item (Sai) Có thể bỏ qua hoàn toàn đơn vị và điều kiện vật lí khi tính toán.
\end{itemchoice}
}
\end{ex}

% ===== Câu 58 =====
% ID: L11C2B10-15-TL
% Mức: VD
\begin{ex}
% ID: L11C2B10-15-TL
% Mức: VD
Trong chuyên đề “Thiết kế phương án thí nghiệm đo tần số của nguồn âm”, Nêu ba lưu ý an toàn và kĩ thuật khi thực hiện thí nghiệm âm.
\choiceTF

\loigiai{
Giữ âm lượng an toàn, nối thiết bị đúng cổng/thang đo, tránh va đập microphone và giảm tiếng ồn nền.
}
\end{ex}

% ===== Câu 59 =====
% ID: L11C2B10-15-TLN
% Mức: TH
\begin{ex}
% ID: L11C2B10-15-TLN
% Mức: TH
Trong chuyên đề “Thiết kế phương án thí nghiệm đo tần số của nguồn âm”, 8 chu kì kéo dài 16 ms. Tần số tính theo Hz bằng bao nhiêu?
\shortans{500}
\loigiai{
Xác định đúng dữ kiện và đơn vị, áp dụng hệ thức của dạng bài rồi tính được kết quả 500.
}
\end{ex}

% ===== Câu 60 =====
% ID: L11C2B10-15-TN
% Mức: TH
\dangbt{Phân tích tín hiệu âm bằng dao động kí và phần mềm}
\begin{ex}
% ID: L11C2B10-15-TN
% Mức: TH
Trong chuyên đề “Phân tích tín hiệu âm bằng dao động kí và phần mềm”, Khoảng thời gian 4,5 ms chứa 3 chu kì. Tần số tính theo Hz, làm tròn đến hàng đơn vị, bằng bao nhiêu? Chọn giá trị đúng.
\choice
{\True 667}
{1334}
{333.5}
{668}
\loigiai{
Đáp án A. 667. Đối chiếu định nghĩa, hệ thức hoặc dữ kiện của bài toán cho kết luận trên.
}
\end{ex}

% ===== Câu 61 =====
% ID: L11C2B10-16-DS
% Mức: TH
\dangbt{Xử lí số liệu nhiều lần đo và báo cáo kết quả tần số}
\begin{ex}
% ID: L11C2B10-16-DS
% Mức: TH
Trong chuyên đề “Xử lí số liệu nhiều lần đo và báo cáo kết quả tần số”, Đánh giá phép đo tần số bằng đồ thị.
\choiceTF
{\True Độ phân giải thời gian ảnh hưởng kết quả.}
{\True Đo lặp giúp phát hiện độ phân tán.}
{\True Có thể lấy trung bình các lần đo hợp lệ.}
{Sai số hệ thống luôn bị loại hết nhờ đo lặp.}
\loigiai{
\begin{itemchoice}
\item (Đúng) Độ phân giải thời gian ảnh hưởng kết quả. (Vì): ộ phân giải thời gian ảnh hưởng kết quả. b) Đ: Đo lặp giúp phát hiện độ phân tán. c) Đ: Có thể lấy trung bình các lần đo hợp lệ. d) S: Sai số hệ thống luôn bị loại hết nhờ đo lặp
\item (Đúng) Đo lặp giúp phát hiện độ phân tán.
\item (Đúng) Có thể lấy trung bình các lần đo hợp lệ.
\item (Sai) Sai số hệ thống luôn bị loại hết nhờ đo lặp.
\end{itemchoice}
}
\end{ex}

% ===== Câu 62 =====
% ID: L11C2B10-16-TL
% Mức: VD
\dangbt{Thiết kế phương án thí nghiệm đo tần số của nguồn âm}
\begin{ex}
% ID: L11C2B10-16-TL
% Mức: VD
Trong chuyên đề “Thiết kế phương án thí nghiệm đo tần số của nguồn âm”, Phân tích các bước lập luận và chỉ ra điều kiện áp dụng trong bài toán: Trình bày cách bố trí thí nghiệm đo tần số của một âm thoa bằng microphone và dao động kí.
\choiceTF

\loigiai{
Đặt microphone gần âm thoa, nối với dao động kí/phần mềm, kích thích âm thoa, điều chỉnh thang thời gian để tín hiệu rõ rồi đo chu kì và tính $f=1/T$. Cần nêu rõ điều kiện áp dụng, đổi đơn vị thống nhất và kiểm tra ý nghĩa vật lí của kết quả.
}
\end{ex}

% ===== Câu 63 =====
% ID: L11C2B10-16-TLN
% Mức: VD
\begin{ex}
% ID: L11C2B10-16-TLN
% Mức: VD
Trong chuyên đề “Thiết kế phương án thí nghiệm đo tần số của nguồn âm”, Từ kết quả của tình huống sau, nếu đại lượng cần tìm được nhân với hệ số 2 thì giá trị mới bằng bao nhiêu? Trên màn hình, 5 chu kì chiếm 10 ms. Chu kì tính theo ms bằng bao nhiêu?
\shortans{4}
\loigiai{
Xác định đúng dữ kiện và đơn vị, áp dụng hệ thức của dạng bài rồi tính được kết quả 4.
}
\end{ex}

% ===== Câu 64 =====
% ID: L11C2B10-16-TN
% Mức: TH
\dangbt{Phân tích tín hiệu âm bằng dao động kí và phần mềm}
\begin{ex}
% ID: L11C2B10-16-TN
% Mức: TH
Trong chuyên đề “Phân tích tín hiệu âm bằng dao động kí và phần mềm”, Hai đỉnh thứ nhất và thứ chín cách nhau 16 ms. Tần số tính theo Hz bằng bao nhiêu? Chọn giá trị đúng.
\choice
{501}
{\True 500}
{1000}
{250}
\loigiai{
Đáp án B. 500. Đối chiếu định nghĩa, hệ thức hoặc dữ kiện của bài toán cho kết luận trên.
}
\end{ex}

% ===== Câu 65 =====
% ID: L11C2B10-17-DS
% Mức: TH
\dangbt{Xử lí số liệu nhiều lần đo và báo cáo kết quả tần số}
\begin{ex}
% ID: L11C2B10-17-DS
% Mức: TH
Trong chuyên đề “Xử lí số liệu nhiều lần đo và báo cáo kết quả tần số”, Một nhóm đo $T=(2,00\pm0,04)$ ms.
\choiceTF
{\True Sai số tương đối của $T$ là 2%.}
{\True Tần số gần bằng $500\,\text{Hz}$.}
{\True Sai số tương đối của f xấp xỉ 2%.}
{Có thể báo f chính xác tuyệt đối $500,000\,\text{Hz}$.}
\loigiai{
\begin{itemchoice}
\item (Đúng) Sai số tương đối của $T$ là 2%. (Vì): Sai số tương đối của $T$ là 2%. b) Đ: Tần số gần bằng $500\,\text{Hz}$. c) Đ: Sai số tương đối của f xấp xỉ 2%. d) S: Có thể báo f chính xác tuyệt đối $500,000\,\text{Hz}$
\item (Đúng) Tần số gần bằng $500\,\text{Hz}$.
\item (Đúng) Sai số tương đối của f xấp xỉ 2%.
\item (Sai) Có thể báo f chính xác tuyệt đối $500,000\,\text{Hz}$.
\end{itemchoice}
}
\end{ex}

% ===== Câu 66 =====
% ID: L11C2B10-17-TL
% Mức: VD
\dangbt{Thiết kế phương án thí nghiệm đo tần số của nguồn âm}
\begin{ex}
% ID: L11C2B10-17-TL
% Mức: VD
Trong chuyên đề “Thiết kế phương án thí nghiệm đo tần số của nguồn âm”, Một học sinh giải bài sau nhưng bỏ qua điều kiện vật lí. Hãy sửa cách làm và trình bày lời giải đúng: Giải thích vì sao nên đo thời gian của nhiều chu kì thay vì chỉ một chu kì.
\choiceTF

\loigiai{
Cách giải đúng: Cùng một độ không đảm bảo khi đọc thời gian, đo khoảng dài hơn làm sai số tương đối nhỏ hơn; sau đó chia cho số chu kì. Sai lầm cần tránh là dùng hệ thức khi chưa xác định điều kiện biên, môi trường hoặc đơn vị.
}
\end{ex}

% ===== Câu 67 =====
% ID: L11C2B10-17-TLN
% Mức: VD
\begin{ex}
% ID: L11C2B10-17-TLN
% Mức: VD
Trong chuyên đề “Thiết kế phương án thí nghiệm đo tần số của nguồn âm”, Từ kết quả của tình huống sau, nếu đại lượng cần tìm được nhân với hệ số 0.5 thì giá trị mới bằng bao nhiêu? Tín hiệu có chu kì 4 ms. Tần số tính theo Hz bằng bao nhiêu?
\shortans{125}
\loigiai{
Xác định đúng dữ kiện và đơn vị, áp dụng hệ thức của dạng bài rồi tính được kết quả 125.
}
\end{ex}

% ===== Câu 68 =====
% ID: L11C2B10-17-TN
% Mức: VD
\dangbt{Phân tích tín hiệu âm bằng dao động kí và phần mềm}
\begin{ex}
% ID: L11C2B10-17-TN
% Mức: VD
Trong chuyên đề “Phân tích tín hiệu âm bằng dao động kí và phần mềm”, Kết luận nào phù hợp nhất khi giải bài toán sau: Đồ thị âm cho thấy từ đỉnh thứ 2 đến đỉnh thứ 12 là 25 ms. Tính chu kì và tần số.
\choice
{Đại lượng cần tìm luôn bằng 0 trong mọi trường hợp.}
{Kết quả không phụ thuộc môi trường, tần số hoặc điều kiện biên.}
{\True Giữa 11 đỉnh có 10 chu kì: $T=25/10=2,5$ ms}
{Không cần sử dụng định nghĩa hay dữ kiện đã cho.}
\loigiai{
Đáp án C. Giữa 11 đỉnh có 10 chu kì: $T=25/10=2,5$ ms. Đối chiếu định nghĩa, hệ thức hoặc dữ kiện của bài toán cho kết luận trên.
}
\end{ex}

% ===== Câu 69 =====
% ID: L11C2B10-18-DS
% Mức: TH
\dangbt{Xử lí số liệu nhiều lần đo và báo cáo kết quả tần số}
\begin{ex}
% ID: L11C2B10-18-DS
% Mức: TH
Trong chuyên đề “Xử lí số liệu nhiều lần đo và báo cáo kết quả tần số”, Khi tín hiệu bị cắt đỉnh do đặt độ nhạy quá cao.
\choiceTF
{\True Dạng sóng bị méo.}
{\True Việc xác định thời điểm đỉnh có thể kém tin cậy.}
{\True Nên điều chỉnh thang biên độ phù hợp.}
{Cắt đỉnh luôn làm chu kì thật của nguồn đổi.}
\loigiai{
\begin{itemchoice}
\item (Đúng) Dạng sóng bị méo. (Vì): ạng sóng bị méo. b) Đ: Việc xác định thời điểm đỉnh có thể kém tin cậy. c) Đ: Nên điều chỉnh thang biên độ phù hợp. d) S: Cắt đỉnh luôn làm chu kì thật của nguồn đổi
\item (Đúng) Việc xác định thời điểm đỉnh có thể kém tin cậy.
\item (Đúng) Nên điều chỉnh thang biên độ phù hợp.
\item (Sai) Cắt đỉnh luôn làm chu kì thật của nguồn đổi.
\end{itemchoice}
}
\end{ex}

% ===== Câu 70 =====
% ID: L11C2B10-18-TL
% Mức: VD
\dangbt{Thiết kế phương án thí nghiệm đo tần số của nguồn âm}
\begin{ex}
% ID: L11C2B10-18-TL
% Mức: VD
Trong chuyên đề “Thiết kế phương án thí nghiệm đo tần số của nguồn âm”, Mở rộng tình huống sau thành một ví dụ thực tế và trình bày phương pháp giải: Nêu ba lưu ý an toàn và kĩ thuật khi thực hiện thí nghiệm âm.
\choiceTF

\loigiai{
Giữ âm lượng an toàn, nối thiết bị đúng cổng/thang đo, tránh va đập microphone và giảm tiếng ồn nền. Có thể kiểm chứng bằng cách thay kết quả vào dữ kiện ban đầu và đối chiếu với hiện tượng thực tế.
}
\end{ex}

% ===== Câu 71 =====
% ID: L11C2B10-18-TLN
% Mức: VD
\begin{ex}
% ID: L11C2B10-18-TLN
% Mức: VD
Trong chuyên đề “Thiết kế phương án thí nghiệm đo tần số của nguồn âm”, Từ kết quả của tình huống sau, nếu đại lượng cần tìm được nhân với hệ số 1.5 thì giá trị mới bằng bao nhiêu? 8 chu kì kéo dài 16 ms. Tần số tính theo Hz bằng bao nhiêu?
\shortans{750}
\loigiai{
Xác định đúng dữ kiện và đơn vị, áp dụng hệ thức của dạng bài rồi tính được kết quả 750.
}
\end{ex}

% ===== Câu 72 =====
% ID: L11C2B10-18-TN
% Mức: VD
\dangbt{Phân tích tín hiệu âm bằng dao động kí và phần mềm}
\begin{ex}
% ID: L11C2B10-18-TN
% Mức: VD
Trong chuyên đề “Phân tích tín hiệu âm bằng dao động kí và phần mềm”, Kết luận nào phù hợp nhất khi giải bài toán sau: Ba lần đo thời gian của 10 chu kì là 19,8 ms; 20,0 ms; 20,2 ms. Tính tần số từ giá trị trung bình.
\choice
{Không cần sử dụng định nghĩa hay dữ kiện đã cho.}
{Đại lượng cần tìm luôn bằng 0 trong mọi trường hợp.}
{Kết quả không phụ thuộc môi trường, tần số hoặc điều kiện biên.}
{\True Thời gian trung bình của 10 chu kì là 20,0 ms}
\loigiai{
Đáp án D. Thời gian trung bình của 10 chu kì là 20,0 ms. Đối chiếu định nghĩa, hệ thức hoặc dữ kiện của bài toán cho kết luận trên.
}
\end{ex}

% ===== Câu 73 =====
% ID: L11C2B10-19-DS
% Mức: VD
\dangbt{Xử lí số liệu nhiều lần đo và báo cáo kết quả tần số}
\begin{ex}
% ID: L11C2B10-19-DS
% Mức: VD
Trong chuyên đề “Xử lí số liệu nhiều lần đo và báo cáo kết quả tần số”, Xét các nhận định tổng hợp thuộc dạng “Xử lí số liệu nhiều lần đo và báo cáo kết quả tần số”.
\choiceTF
{\True xác định vị trí đỉnh không hoàn toàn giống nhau.}
{\True đo nhiều chu kì rồi chia.}
{\True nhỏ hơn giá trị thật.}
{dùng công thức $f=1/T$.}
\loigiai{
\begin{itemchoice}
\item (Đúng) xác định vị trí đỉnh không hoàn toàn giống nhau. (Vì): xác định vị trí đỉnh không hoàn toàn giống nhau. b) Đ: đo nhiều chu kì rồi chia. c) Đ: nhỏ hơn giá trị thật. d) S: dùng công thức $f=1/T$
\item (Đúng) đo nhiều chu kì rồi chia.
\item (Đúng) nhỏ hơn giá trị thật.
\item (Sai) dùng công thức $f=1/T$.
\end{itemchoice}
}
\end{ex}

% ===== Câu 74 =====
% ID: L11C2B10-19-TL
% Mức: TH
\begin{ex}
% ID: L11C2B10-19-TL
% Mức: TH
Trong chuyên đề “Xử lí số liệu nhiều lần đo và báo cáo kết quả tần số”, Phân biệt sai số ngẫu nhiên và sai số hệ thống trong phép đo tần số âm, kèm một ví dụ cho mỗi loại.
\choiceTF

\loigiai{
Sai số ngẫu nhiên thay đổi giữa các lần đọc, ví dụ chọn vị trí đỉnh; sai số hệ thống lệch cùng chiều, ví dụ trục thời gian của thiết bị chưa hiệu chuẩn.
}
\end{ex}

% ===== Câu 75 =====
% ID: L11C2B10-19-TLN
% Mức: TH
\begin{ex}
% ID: L11C2B10-19-TLN
% Mức: TH
Trong chuyên đề “Xử lí số liệu nhiều lần đo và báo cáo kết quả tần số”, Đo được tần số $495\,\text{Hz}$, giá trị chuẩn $500\,\text{Hz}$. Sai số tương đối tính theo phần trăm bằng bao nhiêu?
\shortans{1}
\loigiai{
Xác định đúng dữ kiện và đơn vị, áp dụng hệ thức của dạng bài rồi tính được kết quả 1.
}
\end{ex}

% ===== Câu 76 =====
% ID: L11C2B10-19-TN
% Mức: VD
\dangbt{Phân tích tín hiệu âm bằng dao động kí và phần mềm}
\begin{ex}
% ID: L11C2B10-19-TN
% Mức: VD
Trong chuyên đề “Phân tích tín hiệu âm bằng dao động kí và phần mềm”, Kết luận nào phù hợp nhất khi giải bài toán sau: Mô tả cách nhận biết và loại bỏ một đoạn tín hiệu bị nhiễu trước khi tính tần số.
\choice
{\True Chọn đoạn tuần hoàn ổn định, biên dạng lặp rõ}
{Không cần sử dụng định nghĩa hay dữ kiện đã cho.}
{Đại lượng cần tìm luôn bằng 0 trong mọi trường hợp.}
{Kết quả không phụ thuộc môi trường, tần số hoặc điều kiện biên.}
\loigiai{
Đáp án A. Chọn đoạn tuần hoàn ổn định, biên dạng lặp rõ. Đối chiếu định nghĩa, hệ thức hoặc dữ kiện của bài toán cho kết luận trên.
}
\end{ex}

% ===== Câu 77 =====
% ID: L11C2B10-20-DS
% Mức: VD
\dangbt{Xử lí số liệu nhiều lần đo và báo cáo kết quả tần số}
\begin{ex}
% ID: L11C2B10-20-DS
% Mức: VD
Trong chuyên đề “Xử lí số liệu nhiều lần đo và báo cáo kết quả tần số”, Đánh giá cách xử lí các dữ kiện định lượng của dạng “Xử lí số liệu nhiều lần đo và báo cáo kết quả tần số”.
\choiceTF
{\True Kết quả của bài toán thứ nhất là 1.}
{\True Kết quả của bài toán thứ hai là 500.}
{\True Kết quả của bài toán thứ ba là 5.}
{Có thể bỏ qua hoàn toàn đơn vị và điều kiện vật lí khi tính toán.}
\loigiai{
\begin{itemchoice}
\item (Đúng) Kết quả của bài toán thứ nhất là 1. (Vì): Kết quả của bài toán thứ nhất là 1. b) Đ: Kết quả của bài toán thứ hai là 500. c) Đ: Kết quả của bài toán thứ ba là 5. d) S: Có thể bỏ qua hoàn toàn đơn vị và điều kiện vật lí khi tính toán
\item (Đúng) Kết quả của bài toán thứ hai là 500.
\item (Đúng) Kết quả của bài toán thứ ba là 5.
\item (Sai) Có thể bỏ qua hoàn toàn đơn vị và điều kiện vật lí khi tính toán.
\end{itemchoice}
}
\end{ex}

% ===== Câu 78 =====
% ID: L11C2B10-20-TL
% Mức: TH
\begin{ex}
% ID: L11C2B10-20-TL
% Mức: TH
Trong chuyên đề “Xử lí số liệu nhiều lần đo và báo cáo kết quả tần số”, Một nhóm đo $T=(2,00\pm0,04)$ ms. Ước lượng tần số và độ không đảm bảo tuyệt đối của tần số.
\choiceTF

\loigiai{
$f=500$ Hz; sai số tương đối xấp xỉ 2%, nên $\Delta f\approx10$ Hz. Kết quả $f=(500\pm10)$ Hz.
}
\end{ex}

% ===== Câu 79 =====
% ID: L11C2B10-20-TLN
% Mức: TH
\begin{ex}
% ID: L11C2B10-20-TLN
% Mức: TH
Trong chuyên đề “Xử lí số liệu nhiều lần đo và báo cáo kết quả tần số”, Ba kết quả $498\,\text{Hz}$, $500\,\text{Hz}$, $502\,\text{Hz}$. Giá trị trung bình tính theo Hz bằng bao nhiêu?
\shortans{14.5}
\loigiai{
Xác định đúng dữ kiện và đơn vị, áp dụng hệ thức của dạng bài rồi tính được kết quả 500.
}
\end{ex}

% ===== Câu 80 =====
% ID: L11C2B10-20-TN
% Mức: VD
\dangbt{Phân tích tín hiệu âm bằng dao động kí và phần mềm}
\begin{ex}
% ID: L11C2B10-20-TN
% Mức: VD
Trong chuyên đề “Phân tích tín hiệu âm bằng dao động kí và phần mềm”, Phương pháp phù hợp nhất cho dạng “Phân tích tín hiệu âm bằng dao động kí và phần mềm” là
\choice
{luôn coi mọi điểm dao động cùng pha}
{\True xác định dữ kiện, chọn đúng hệ thức hoặc định nghĩa, đổi đơn vị rồi kiểm tra kết quả}
{chỉ thay số mà không cần xét điều kiện áp dụng}
{bỏ qua đơn vị và chiều truyền sóng}
\loigiai{
Đáp án B. xác định dữ kiện, chọn đúng hệ thức hoặc định nghĩa, đổi đơn vị rồi kiểm tra kết quả. Đối chiếu định nghĩa, hệ thức hoặc dữ kiện của bài toán cho kết luận trên.
}
\end{ex}

% ===== Câu 81 =====
% ID: L11C2B10-21-DS
% Mức: TH
\dangbt{Xử lí đồ thị dao động âm để xác định chu kì và tần số}
\begin{ex}
% ID: L11C2B10-21-DS
% Mức: TH
Đồ thị cho 6 chu kì trong 12 ms.
\choiceTF
{\True Chu kì là 2 ms.}
{\True Tần số là $500\,\text{Hz}$.}
{\True Trong $1\,\text{s}$ có 500 dao động.}
{Tăng biên độ hiển thị làm tần số tăng.}
\loigiai{
\begin{itemchoice}
\item (Đúng) Chu kì là 2 ms. (Vì): Chu kì là 2 ms. b) Đ: Tần số là $500\,\text{Hz}$. c) Đ: Trong $1\,\text{s}$ có 500 dao động. d) S: Tăng biên độ hiển thị làm tần số tăng
\item (Đúng) Tần số là $500\,\text{Hz}$.
\item (Đúng) Trong $1\,\text{s}$ có 500 dao động.
\item (Sai) Tăng biên độ hiển thị làm tần số tăng.
\end{itemchoice}
}
\end{ex}

% ===== Câu 82 =====
% ID: L11C2B10-21-TL
% Mức: VD
\dangbt{Xử lí số liệu nhiều lần đo và báo cáo kết quả tần số}
\begin{ex}
% ID: L11C2B10-21-TL
% Mức: VD
Trong chuyên đề “Xử lí số liệu nhiều lần đo và báo cáo kết quả tần số”, Đề xuất ba cải tiến để tăng độ tin cậy của phép đo tần số âm.
\choiceTF

\loigiai{
Giảm tiếng ồn, đo nhiều chu kì và nhiều lần, chọn thang thời gian phù hợp/hiệu chuẩn thiết bị, giữ nguồn và microphone ổn định.
}
\end{ex}

% ===== Câu 83 =====
% ID: L11C2B10-21-TLN
% Mức: TH
\begin{ex}
% ID: L11C2B10-21-TLN
% Mức: TH
Trong chuyên đề “Xử lí số liệu nhiều lần đo và báo cáo kết quả tần số”, Chu kì đo được 2,0 ms với sai số 0,1 ms. Sai số tương đối của chu kì tính theo phần trăm bằng bao nhiêu?
\shortans{5}
\loigiai{
Xác định đúng dữ kiện và đơn vị, áp dụng hệ thức của dạng bài rồi tính được kết quả 5.
}
\end{ex}

% ===== Câu 84 =====
% ID: L11C2B10-21-TN
% Mức: NB
\dangbt{Thiết kế phương án thí nghiệm đo tần số của nguồn âm}
\begin{ex}
% ID: L11C2B10-21-TN
% Mức: NB
Trong chuyên đề “Thiết kế phương án thí nghiệm đo tần số của nguồn âm”, Thiết bị dùng để thu âm và biến dao động âm thành tín hiệu điện là
\choice
{thước mét}
{nhiệt kế}
{lực kế}
{\True microphone}
\loigiai{
Đáp án D. microphone. Đối chiếu định nghĩa, hệ thức hoặc dữ kiện của bài toán cho kết luận trên.
}
\end{ex}

% ===== Câu 85 =====
% ID: L11C2B10-22-DS
% Mức: TH
\dangbt{Xử lí đồ thị dao động âm để xác định chu kì và tần số}
\begin{ex}
% ID: L11C2B10-22-DS
% Mức: TH
Hai lần đo cho 5 chu kì lần lượt là 10,0 ms và 10,2 ms.
\choiceTF
{\True Có thể lấy trung bình thời gian trước khi tính chu kì.}
{\True Thời gian trung bình của 5 chu kì là 10,1 ms.}
{\True Chu kì trung bình là 2,02 ms.}
{Tần số trung bình chính xác bằng $202\,\text{Hz}$.}
\loigiai{
\begin{itemchoice}
\item (Đúng) Có thể lấy trung bình thời gian trước khi tính chu kì. (Vì): Có thể lấy trung bình thời gian trước khi tính chu kì. b) Đ: Thời gian trung bình của 5 chu kì là 10,1 ms. c) Đ: Chu kì trung bình là 2,02 ms. d) S: Tần số trung bình chính xác bằng $202\,\text{Hz}$
\item (Đúng) Thời gian trung bình của 5 chu kì là 10,1 ms.
\item (Đúng) Chu kì trung bình là 2,02 ms.
\item (Sai) Tần số trung bình chính xác bằng $202\,\text{Hz}$.
\end{itemchoice}
}
\end{ex}

% ===== Câu 86 =====
% ID: L11C2B10-22-TL
% Mức: VD
\dangbt{Xử lí số liệu nhiều lần đo và báo cáo kết quả tần số}
\begin{ex}
% ID: L11C2B10-22-TL
% Mức: VD
Trong chuyên đề “Xử lí số liệu nhiều lần đo và báo cáo kết quả tần số”, Phân tích các bước lập luận và chỉ ra điều kiện áp dụng trong bài toán: Phân biệt sai số ngẫu nhiên và sai số hệ thống trong phép đo tần số âm, kèm một ví dụ cho mỗi loại.
\choiceTF

\loigiai{
Sai số ngẫu nhiên thay đổi giữa các lần đọc, ví dụ chọn vị trí đỉnh; sai số hệ thống lệch cùng chiều, ví dụ trục thời gian của thiết bị chưa hiệu chuẩn. Cần nêu rõ điều kiện áp dụng, đổi đơn vị thống nhất và kiểm tra ý nghĩa vật lí của kết quả.
}
\end{ex}

% ===== Câu 87 =====
% ID: L11C2B10-22-TLN
% Mức: VD
\begin{ex}
% ID: L11C2B10-22-TLN
% Mức: VD
Trong chuyên đề “Xử lí số liệu nhiều lần đo và báo cáo kết quả tần số”, Từ kết quả của tình huống sau, nếu đại lượng cần tìm được nhân với hệ số 2 thì giá trị mới bằng bao nhiêu? Đo được tần số $495\,\text{Hz}$, giá trị chuẩn $500\,\text{Hz}$. Sai số tương đối tính theo phần trăm bằng bao nhiêu?
\shortans{2}
\loigiai{
Xác định đúng dữ kiện và đơn vị, áp dụng hệ thức của dạng bài rồi tính được kết quả 2.
}
\end{ex}

% ===== Câu 88 =====
% ID: L11C2B10-22-TN
% Mức: NB
\dangbt{Thiết kế phương án thí nghiệm đo tần số của nguồn âm}
\begin{ex}
% ID: L11C2B10-22-TN
% Mức: NB
Trong chuyên đề “Thiết kế phương án thí nghiệm đo tần số của nguồn âm”, Để quan sát dạng tín hiệu âm theo thời gian cần dùng
\choice
{\True dao động kí hoặc phần mềm ghi âm}
{cân điện tử}
{ampe kế một chiều}
{la bàn}
\loigiai{
Đáp án A. dao động kí hoặc phần mềm ghi âm. Đối chiếu định nghĩa, hệ thức hoặc dữ kiện của bài toán cho kết luận trên.
}
\end{ex}

% ===== Câu 89 =====
% ID: L11C2B10-23-DS
% Mức: TH
\dangbt{Xử lí đồ thị dao động âm để xác định chu kì và tần số}
\begin{ex}
% ID: L11C2B10-23-DS
% Mức: TH
Một tín hiệu âm có chu kì 1,25 ms.
\choiceTF
{\True Chu kì bằng $0,00125\,\text{s}$.}
{\True Tần số bằng $800\,\text{Hz}$.}
{\True Nếu màn hình hiển thị 4 chu kì thì khoảng thời gian là 5 ms.}
{Âm này là siêu âm.}
\loigiai{
\begin{itemchoice}
\item (Đúng) Chu kì bằng $0,00125\,\text{s}$. (Vì): Chu kì bằng $0,00125\,\text{s}$. b) Đ: Tần số bằng $800\,\text{Hz}$. c) Đ: Nếu màn hình hiển thị 4 chu kì thì khoảng thời gian là 5 ms. d) S: Âm này là siêu âm
\item (Đúng) Tần số bằng $800\,\text{Hz}$.
\item (Đúng) Nếu màn hình hiển thị 4 chu kì thì khoảng thời gian là 5 ms.
\item (Sai) Âm này là siêu âm.
\end{itemchoice}
}
\end{ex}

% ===== Câu 90 =====
% ID: L11C2B10-23-TL
% Mức: VD
\dangbt{Xử lí số liệu nhiều lần đo và báo cáo kết quả tần số}
\begin{ex}
% ID: L11C2B10-23-TL
% Mức: VD
Trong chuyên đề “Xử lí số liệu nhiều lần đo và báo cáo kết quả tần số”, Một học sinh giải bài sau nhưng bỏ qua điều kiện vật lí. Hãy sửa cách làm và trình bày lời giải đúng: Một nhóm đo $T=(2,00\pm0,04)$ ms. Ước lượng tần số và độ không đảm bảo tuyệt đối của tần số.
\choiceTF

\loigiai{
Cách giải đúng: $f=500$ Hz; sai số tương đối xấp xỉ 2%, nên $\Delta f\approx10$ Hz. Kết quả $f=(500\pm10)$ Hz. Sai lầm cần tránh là dùng hệ thức khi chưa xác định điều kiện biên, môi trường hoặc đơn vị.
}
\end{ex}

% ===== Câu 91 =====
% ID: L11C2B10-23-TLN
% Mức: VD
\begin{ex}
% ID: L11C2B10-23-TLN
% Mức: VD
Trong chuyên đề “Xử lí số liệu nhiều lần đo và báo cáo kết quả tần số”, Từ kết quả của tình huống sau, nếu đại lượng cần tìm được nhân với hệ số 0.5 thì giá trị mới bằng bao nhiêu? Ba kết quả $498\,\text{Hz}$, $500\,\text{Hz}$, $502\,\text{Hz}$. Giá trị trung bình tính theo Hz bằng bao nhiêu?
\shortans{250}
\loigiai{
Xác định đúng dữ kiện và đơn vị, áp dụng hệ thức của dạng bài rồi tính được kết quả 250.
}
\end{ex}

% ===== Câu 92 =====
% ID: L11C2B10-23-TN
% Mức: TH
\dangbt{Thiết kế phương án thí nghiệm đo tần số của nguồn âm}
\begin{ex}
% ID: L11C2B10-23-TN
% Mức: TH
Trong chuyên đề “Thiết kế phương án thí nghiệm đo tần số của nguồn âm”, Khi đo tần số âm, đại lượng cần đọc trực tiếp trên đồ thị là
\choice
{cường độ dòng điện}
{\True chu kì hoặc thời gian của nhiều chu kì}
{bước sóng trong chân không}
{khối lượng nguồn}
\loigiai{
Đáp án B. chu kì hoặc thời gian của nhiều chu kì. Đối chiếu định nghĩa, hệ thức hoặc dữ kiện của bài toán cho kết luận trên.
}
\end{ex}

% ===== Câu 93 =====
% ID: L11C2B10-24-DS
% Mức: VD
\dangbt{Xử lí đồ thị dao động âm để xác định chu kì và tần số}
\begin{ex}
% ID: L11C2B10-24-DS
% Mức: VD
Xét các nhận định tổng hợp thuộc dạng “Xử lí đồ thị dao động âm để xác định chu kì và tần số”.
\choiceTF
{\True $500\,\text{Hz}$.}
{\True $400\,\text{Hz}$.}
{\True đo thời gian của nhiều chu kì.}
{$50\,\text{Hz}$.}
\loigiai{
\begin{itemchoice}
\item (Đúng) $500\,\text{Hz}$. (Vì): $500\,\text{Hz}$. b) Đ: $400\,\text{Hz}$. c) Đ: đo thời gian của nhiều chu kì. d) S: $50\,\text{Hz}$
\item (Đúng) $400\,\text{Hz}$.
\item (Đúng) đo thời gian của nhiều chu kì.
\item (Sai) $50\,\text{Hz}$.
\end{itemchoice}
}
\end{ex}

% ===== Câu 94 =====
% ID: L11C2B10-24-TL
% Mức: VD
\dangbt{Xử lí số liệu nhiều lần đo và báo cáo kết quả tần số}
\begin{ex}
% ID: L11C2B10-24-TL
% Mức: VD
Trong chuyên đề “Xử lí số liệu nhiều lần đo và báo cáo kết quả tần số”, Mở rộng tình huống sau thành một ví dụ thực tế và trình bày phương pháp giải: Đề xuất ba cải tiến để tăng độ tin cậy của phép đo tần số âm.
\choiceTF

\loigiai{
Giảm tiếng ồn, đo nhiều chu kì và nhiều lần, chọn thang thời gian phù hợp/hiệu chuẩn thiết bị, giữ nguồn và microphone ổn định. Có thể kiểm chứng bằng cách thay kết quả vào dữ kiện ban đầu và đối chiếu với hiện tượng thực tế.
}
\end{ex}

% ===== Câu 95 =====
% ID: L11C2B10-24-TLN
% Mức: VD
\begin{ex}
% ID: L11C2B10-24-TLN
% Mức: VD
Trong chuyên đề “Xử lí số liệu nhiều lần đo và báo cáo kết quả tần số”, Từ kết quả của tình huống sau, nếu đại lượng cần tìm được nhân với hệ số 1.5 thì giá trị mới bằng bao nhiêu? Chu kì đo được 2,0 ms với sai số 0,1 ms. Sai số tương đối của chu kì tính theo phần trăm bằng bao nhiêu?
\shortans{07/05/2026}
\loigiai{
Xác định đúng dữ kiện và đơn vị, áp dụng hệ thức của dạng bài rồi tính được kết quả 7.5.
}
\end{ex}

% ===== Câu 96 =====
% ID: L11C2B10-24-TN
% Mức: TH
\dangbt{Thiết kế phương án thí nghiệm đo tần số của nguồn âm}
\begin{ex}
% ID: L11C2B10-24-TN
% Mức: TH
Trong chuyên đề “Thiết kế phương án thí nghiệm đo tần số của nguồn âm”, Trên màn hình, 5 chu kì chiếm 10 ms. Chu kì tính theo ms bằng bao nhiêu? Chọn giá trị đúng.
\choice
{1}
{3}
{\True 2}
{4}
\loigiai{
Đáp án C. 2. Đối chiếu định nghĩa, hệ thức hoặc dữ kiện của bài toán cho kết luận trên.
}
\end{ex}

% ===== Câu 97 =====
% ID: L11C2B10-25-DS
% Mức: VD
\dangbt{Xử lí đồ thị dao động âm để xác định chu kì và tần số}
\begin{ex}
% ID: L11C2B10-25-DS
% Mức: VD
Đánh giá cách xử lí các dữ kiện định lượng của dạng “Xử lí đồ thị dao động âm để xác định chu kì và tần số”.
\choiceTF
{\True Kết quả của bài toán thứ nhất là 500.}
{\True Kết quả của bài toán thứ hai là 667.}
{\True Kết quả của bài toán thứ ba là 500.}
{Có thể bỏ qua hoàn toàn đơn vị và điều kiện vật lí khi tính toán.}
\loigiai{
\begin{itemchoice}
\item (Đúng) Kết quả của bài toán thứ nhất là 500. (Vì): Kết quả của bài toán thứ nhất là 500. b) Đ: Kết quả của bài toán thứ hai là 667. c) Đ: Kết quả của bài toán thứ ba là 500. d) S: Có thể bỏ qua hoàn toàn đơn vị và điều kiện vật lí khi tính toán
\item (Đúng) Kết quả của bài toán thứ hai là 667.
\item (Đúng) Kết quả của bài toán thứ ba là 500.
\item (Sai) Có thể bỏ qua hoàn toàn đơn vị và điều kiện vật lí khi tính toán.
\end{itemchoice}
}
\end{ex}

% ===== Câu 98 =====
% ID: L11C2B10-25-TL
% Mức: TH
\begin{ex}
% ID: L11C2B10-25-TL
% Mức: TH
Đồ thị âm cho thấy từ đỉnh thứ 2 đến đỉnh thứ 12 là 25 ms. Tính chu kì và tần số.
\choiceTF

\loigiai{
Giữa 11 đỉnh có 10 chu kì: $T=25/10=2,5$ ms; $f=1/(2,5\times10^{-3})=400$ Hz.
}
\end{ex}

% ===== Câu 99 =====
% ID: L11C2B10-25-TLN
% Mức: TH
\begin{ex}
% ID: L11C2B10-25-TLN
% Mức: TH
12 chu kì kéo dài 24 ms. Tần số tính theo Hz bằng bao nhiêu?
\shortans{500}
\loigiai{
Xác định đúng dữ kiện và đơn vị, áp dụng hệ thức của dạng bài rồi tính được kết quả 500.
}
\end{ex}

% ===== Câu 100 =====
% ID: L11C2B10-25-TN
% Mức: TH
\dangbt{Thiết kế phương án thí nghiệm đo tần số của nguồn âm}
\begin{ex}
% ID: L11C2B10-25-TN
% Mức: TH
Trong chuyên đề “Thiết kế phương án thí nghiệm đo tần số của nguồn âm”, Tín hiệu có chu kì 4 ms. Tần số tính theo Hz bằng bao nhiêu? Chọn giá trị đúng.
\choice
{500}
{125}
{251}
{\True 250}
\loigiai{
Đáp án D. 250. Đối chiếu định nghĩa, hệ thức hoặc dữ kiện của bài toán cho kết luận trên.
}
\end{ex}

% ===== Câu 101 =====
% ID: L11C2B10-26-DS
% Mức: TH
\dangbt{Đánh giá kết quả và sai số phép đo tần số âm}
\begin{ex}
% ID: L11C2B10-26-DS
% Mức: TH
Đánh giá phép đo tần số bằng đồ thị.
\choiceTF
{\True Độ phân giải thời gian ảnh hưởng kết quả.}
{\True Đo lặp giúp phát hiện độ phân tán.}
{\True Có thể lấy trung bình các lần đo hợp lệ.}
{Sai số hệ thống luôn bị loại hết nhờ đo lặp.}
\loigiai{
\begin{itemchoice}
\item (Đúng) Độ phân giải thời gian ảnh hưởng kết quả. (Vì): ộ phân giải thời gian ảnh hưởng kết quả. b) Đ: Đo lặp giúp phát hiện độ phân tán. c) Đ: Có thể lấy trung bình các lần đo hợp lệ. d) S: Sai số hệ thống luôn bị loại hết nhờ đo lặp
\item (Đúng) Đo lặp giúp phát hiện độ phân tán.
\item (Đúng) Có thể lấy trung bình các lần đo hợp lệ.
\item (Sai) Sai số hệ thống luôn bị loại hết nhờ đo lặp.
\end{itemchoice}
}
\end{ex}

% ===== Câu 102 =====
% ID: L11C2B10-26-TL
% Mức: TH
\dangbt{Xử lí đồ thị dao động âm để xác định chu kì và tần số}
\begin{ex}
% ID: L11C2B10-26-TL
% Mức: TH
Ba lần đo thời gian của 10 chu kì là 19,8 ms; 20,0 ms; 20,2 ms. Tính tần số từ giá trị trung bình.
\choiceTF

\loigiai{
Thời gian trung bình của 10 chu kì là 20,0 ms; $T=2,00$ ms; $f=500$ Hz.
}
\end{ex}

% ===== Câu 103 =====
% ID: L11C2B10-26-TLN
% Mức: TH
\begin{ex}
% ID: L11C2B10-26-TLN
% Mức: TH
Khoảng thời gian 4,5 ms chứa 3 chu kì. Tần số tính theo Hz, làm tròn đến hàng đơn vị, bằng bao nhiêu?
\shortans{667}
\loigiai{
Xác định đúng dữ kiện và đơn vị, áp dụng hệ thức của dạng bài rồi tính được kết quả 667.
}
\end{ex}

% ===== Câu 104 =====
% ID: L11C2B10-26-TN
% Mức: TH
\dangbt{Thiết kế phương án thí nghiệm đo tần số của nguồn âm}
\begin{ex}
% ID: L11C2B10-26-TN
% Mức: TH
Trong chuyên đề “Thiết kế phương án thí nghiệm đo tần số của nguồn âm”, 8 chu kì kéo dài 16 ms. Tần số tính theo Hz bằng bao nhiêu? Chọn giá trị đúng.
\choice
{\True 500}
{1000}
{250}
{501}
\loigiai{
Đáp án A. 500. Đối chiếu định nghĩa, hệ thức hoặc dữ kiện của bài toán cho kết luận trên.
}
\end{ex}

% ===== Câu 105 =====
% ID: L11C2B10-27-DS
% Mức: TH
\dangbt{Đánh giá kết quả và sai số phép đo tần số âm}
\begin{ex}
% ID: L11C2B10-27-DS
% Mức: TH
Một nhóm đo $T=(2,00\pm0,04)$ ms.
\choiceTF
{\True Sai số tương đối của $T$ là 2%.}
{\True Tần số gần bằng $500\,\text{Hz}$.}
{\True Sai số tương đối của f xấp xỉ 2%.}
{Có thể báo f chính xác tuyệt đối $500,000\,\text{Hz}$.}
\loigiai{
\begin{itemchoice}
\item (Đúng) Sai số tương đối của $T$ là 2%. (Vì): Sai số tương đối của $T$ là 2%. b) Đ: Tần số gần bằng $500\,\text{Hz}$. c) Đ: Sai số tương đối của f xấp xỉ 2%. d) S: Có thể báo f chính xác tuyệt đối $500,000\,\text{Hz}$
\item (Đúng) Tần số gần bằng $500\,\text{Hz}$.
\item (Đúng) Sai số tương đối của f xấp xỉ 2%.
\item (Sai) Có thể báo f chính xác tuyệt đối $500,000\,\text{Hz}$.
\end{itemchoice}
}
\end{ex}

% ===== Câu 106 =====
% ID: L11C2B10-27-TL
% Mức: VD
\dangbt{Xử lí đồ thị dao động âm để xác định chu kì và tần số}
\begin{ex}
% ID: L11C2B10-27-TL
% Mức: VD
Mô tả cách nhận biết và loại bỏ một đoạn tín hiệu bị nhiễu trước khi tính tần số.
\choiceTF

\loigiai{
Chọn đoạn tuần hoàn ổn định, biên dạng lặp rõ; bỏ đoạn có xung bất thường hoặc mất chu kì, rồi đo nhiều chu kì trên phần còn lại.
}
\end{ex}

% ===== Câu 107 =====
% ID: L11C2B10-27-TLN
% Mức: TH
\begin{ex}
% ID: L11C2B10-27-TLN
% Mức: TH
Hai đỉnh thứ nhất và thứ chín cách nhau 16 ms. Tần số tính theo Hz bằng bao nhiêu?
\shortans{500}
\loigiai{
Xác định đúng dữ kiện và đơn vị, áp dụng hệ thức của dạng bài rồi tính được kết quả 500.
}
\end{ex}

% ===== Câu 108 =====
% ID: L11C2B10-27-TN
% Mức: VD
\dangbt{Thiết kế phương án thí nghiệm đo tần số của nguồn âm}
\begin{ex}
% ID: L11C2B10-27-TN
% Mức: VD
Trong chuyên đề “Thiết kế phương án thí nghiệm đo tần số của nguồn âm”, Kết luận nào phù hợp nhất khi giải bài toán sau: Trình bày cách bố trí thí nghiệm đo tần số của một âm thoa bằng microphone và dao động kí.
\choice
{Kết quả không phụ thuộc môi trường, tần số hoặc điều kiện biên.}
{\True Đặt microphone gần âm thoa, nối với dao động kí/phần mềm, kích thích âm thoa, điều chỉnh thang thời gian để tín hiệu rõ rồi đo chu kì và tính $f=1/T$.}
{Không cần sử dụng định nghĩa hay dữ kiện đã cho.}
{Đại lượng cần tìm luôn bằng 0 trong mọi trường hợp.}
\loigiai{
Đáp án B. Đặt microphone gần âm thoa, nối với dao động kí/phần mềm, kích thích âm thoa, điều chỉnh thang thời gian để tín hiệu rõ rồi đo chu kì và tính $f=1/T$.. Đối chiếu định nghĩa, hệ thức hoặc dữ kiện của bài toán cho kết luận trên.
}
\end{ex}

% ===== Câu 109 =====
% ID: L11C2B10-28-DS
% Mức: TH
\dangbt{Đánh giá kết quả và sai số phép đo tần số âm}
\begin{ex}
% ID: L11C2B10-28-DS
% Mức: TH
Khi tín hiệu bị cắt đỉnh do đặt độ nhạy quá cao.
\choiceTF
{\True Dạng sóng bị méo.}
{\True Việc xác định thời điểm đỉnh có thể kém tin cậy.}
{\True Nên điều chỉnh thang biên độ phù hợp.}
{Cắt đỉnh luôn làm chu kì thật của nguồn đổi.}
\loigiai{
\begin{itemchoice}
\item (Đúng) Dạng sóng bị méo. (Vì): ạng sóng bị méo. b) Đ: Việc xác định thời điểm đỉnh có thể kém tin cậy. c) Đ: Nên điều chỉnh thang biên độ phù hợp. d) S: Cắt đỉnh luôn làm chu kì thật của nguồn đổi
\item (Đúng) Việc xác định thời điểm đỉnh có thể kém tin cậy.
\item (Đúng) Nên điều chỉnh thang biên độ phù hợp.
\item (Sai) Cắt đỉnh luôn làm chu kì thật của nguồn đổi.
\end{itemchoice}
}
\end{ex}

% ===== Câu 110 =====
% ID: L11C2B10-28-TL
% Mức: VD
\dangbt{Xử lí đồ thị dao động âm để xác định chu kì và tần số}
\begin{ex}
% ID: L11C2B10-28-TL
% Mức: VD
Phân tích các bước lập luận và chỉ ra điều kiện áp dụng trong bài toán: Đồ thị âm cho thấy từ đỉnh thứ 2 đến đỉnh thứ 12 là 25 ms. Tính chu kì và tần số.
\choiceTF

\loigiai{
Giữa 11 đỉnh có 10 chu kì: $T=25/10=2,5$ ms; $f=1/(2,5\times10^{-3})=400$ Hz. Cần nêu rõ điều kiện áp dụng, đổi đơn vị thống nhất và kiểm tra ý nghĩa vật lí của kết quả.
}
\end{ex}

% ===== Câu 111 =====
% ID: L11C2B10-28-TLN
% Mức: VD
\begin{ex}
% ID: L11C2B10-28-TLN
% Mức: VD
Từ kết quả của tình huống sau, nếu đại lượng cần tìm được nhân với hệ số 2 thì giá trị mới bằng bao nhiêu? 12 chu kì kéo dài 24 ms. Tần số tính theo Hz bằng bao nhiêu?
\shortans{1000}
\loigiai{
Xác định đúng dữ kiện và đơn vị, áp dụng hệ thức của dạng bài rồi tính được kết quả 1000.
}
\end{ex}

% ===== Câu 112 =====
% ID: L11C2B10-28-TN
% Mức: VD
\dangbt{Thiết kế phương án thí nghiệm đo tần số của nguồn âm}
\begin{ex}
% ID: L11C2B10-28-TN
% Mức: VD
Trong chuyên đề “Thiết kế phương án thí nghiệm đo tần số của nguồn âm”, Kết luận nào phù hợp nhất khi giải bài toán sau: Giải thích vì sao nên đo thời gian của nhiều chu kì thay vì chỉ một chu kì.
\choice
{Đại lượng cần tìm luôn bằng 0 trong mọi trường hợp.}
{Kết quả không phụ thuộc môi trường, tần số hoặc điều kiện biên.}
{\True Cùng một độ không đảm bảo khi đọc thời gian, đo khoảng dài hơn làm sai số tương đối nhỏ hơn}
{Không cần sử dụng định nghĩa hay dữ kiện đã cho.}
\loigiai{
Đáp án C. Cùng một độ không đảm bảo khi đọc thời gian, đo khoảng dài hơn làm sai số tương đối nhỏ hơn. Đối chiếu định nghĩa, hệ thức hoặc dữ kiện của bài toán cho kết luận trên.
}
\end{ex}

% ===== Câu 113 =====
% ID: L11C2B10-29-DS
% Mức: VD
\dangbt{Đánh giá kết quả và sai số phép đo tần số âm}
\begin{ex}
% ID: L11C2B10-29-DS
% Mức: VD
Xét các nhận định tổng hợp thuộc dạng “Đánh giá kết quả và sai số phép đo tần số âm”.
\choiceTF
{\True xác định vị trí đỉnh không hoàn toàn giống nhau.}
{\True đo nhiều chu kì rồi chia.}
{\True nhỏ hơn giá trị thật.}
{dùng công thức $f=1/T$.}
\loigiai{
\begin{itemchoice}
\item (Đúng) xác định vị trí đỉnh không hoàn toàn giống nhau. (Vì): xác định vị trí đỉnh không hoàn toàn giống nhau. b) Đ: đo nhiều chu kì rồi chia. c) Đ: nhỏ hơn giá trị thật. d) S: dùng công thức $f=1/T$
\item (Đúng) đo nhiều chu kì rồi chia.
\item (Đúng) nhỏ hơn giá trị thật.
\item (Sai) dùng công thức $f=1/T$.
\end{itemchoice}
}
\end{ex}

% ===== Câu 114 =====
% ID: L11C2B10-29-TL
% Mức: VD
\dangbt{Xử lí đồ thị dao động âm để xác định chu kì và tần số}
\begin{ex}
% ID: L11C2B10-29-TL
% Mức: VD
Một học sinh giải bài sau nhưng bỏ qua điều kiện vật lí. Hãy sửa cách làm và trình bày lời giải đúng: Ba lần đo thời gian của 10 chu kì là 19,8 ms; 20,0 ms; 20,2 ms. Tính tần số từ giá trị trung bình.
\choiceTF

\loigiai{
Cách giải đúng: Thời gian trung bình của 10 chu kì là 20,0 ms; $T=2,00$ ms; $f=500$ Hz. Sai lầm cần tránh là dùng hệ thức khi chưa xác định điều kiện biên, môi trường hoặc đơn vị.
}
\end{ex}

% ===== Câu 115 =====
% ID: L11C2B10-29-TLN
% Mức: VD
\begin{ex}
% ID: L11C2B10-29-TLN
% Mức: VD
Từ kết quả của tình huống sau, nếu đại lượng cần tìm được nhân với hệ số 0.5 thì giá trị mới bằng bao nhiêu? Khoảng thời gian 4,5 ms chứa 3 chu kì. Tần số tính theo Hz, làm tròn đến hàng đơn vị, bằng bao nhiêu?
\shortans{333.5}
\loigiai{
Xác định đúng dữ kiện và đơn vị, áp dụng hệ thức của dạng bài rồi tính được kết quả 333.5.
}
\end{ex}

% ===== Câu 116 =====
% ID: L11C2B10-29-TN
% Mức: VD
\dangbt{Thiết kế phương án thí nghiệm đo tần số của nguồn âm}
\begin{ex}
% ID: L11C2B10-29-TN
% Mức: VD
Trong chuyên đề “Thiết kế phương án thí nghiệm đo tần số của nguồn âm”, Kết luận nào phù hợp nhất khi giải bài toán sau: Nêu ba lưu ý an toàn và kĩ thuật khi thực hiện thí nghiệm âm.
\choice
{Không cần sử dụng định nghĩa hay dữ kiện đã cho.}
{Đại lượng cần tìm luôn bằng 0 trong mọi trường hợp.}
{Kết quả không phụ thuộc môi trường, tần số hoặc điều kiện biên.}
{\True Giữ âm lượng an toàn, nối thiết bị đúng cổng/thang đo, tránh va đập microphone và giảm tiếng ồn nền.}
\loigiai{
Đáp án D. Giữ âm lượng an toàn, nối thiết bị đúng cổng/thang đo, tránh va đập microphone và giảm tiếng ồn nền.. Đối chiếu định nghĩa, hệ thức hoặc dữ kiện của bài toán cho kết luận trên.
}
\end{ex}

% ===== Câu 117 =====
% ID: L11C2B10-30-DS
% Mức: VD
\dangbt{Đánh giá kết quả và sai số phép đo tần số âm}
\begin{ex}
% ID: L11C2B10-30-DS
% Mức: VD
Đánh giá cách xử lí các dữ kiện định lượng của dạng “Đánh giá kết quả và sai số phép đo tần số âm”.
\choiceTF
{\True Kết quả của bài toán thứ nhất là 1.}
{\True Kết quả của bài toán thứ hai là 500.}
{\True Kết quả của bài toán thứ ba là 5.}
{Có thể bỏ qua hoàn toàn đơn vị và điều kiện vật lí khi tính toán.}
\loigiai{
\begin{itemchoice}
\item (Đúng) Kết quả của bài toán thứ nhất là 1. (Vì): Kết quả của bài toán thứ nhất là 1. b) Đ: Kết quả của bài toán thứ hai là 500. c) Đ: Kết quả của bài toán thứ ba là 5. d) S: Có thể bỏ qua hoàn toàn đơn vị và điều kiện vật lí khi tính toán
\item (Đúng) Kết quả của bài toán thứ hai là 500.
\item (Đúng) Kết quả của bài toán thứ ba là 5.
\item (Sai) Có thể bỏ qua hoàn toàn đơn vị và điều kiện vật lí khi tính toán.
\end{itemchoice}
}
\end{ex}

% ===== Câu 118 =====
% ID: L11C2B10-30-TL
% Mức: VD
\dangbt{Xử lí đồ thị dao động âm để xác định chu kì và tần số}
\begin{ex}
% ID: L11C2B10-30-TL
% Mức: VD
Mở rộng tình huống sau thành một ví dụ thực tế và trình bày phương pháp giải: Mô tả cách nhận biết và loại bỏ một đoạn tín hiệu bị nhiễu trước khi tính tần số.
\choiceTF

\loigiai{
Chọn đoạn tuần hoàn ổn định, biên dạng lặp rõ; bỏ đoạn có xung bất thường hoặc mất chu kì, rồi đo nhiều chu kì trên phần còn lại. Có thể kiểm chứng bằng cách thay kết quả vào dữ kiện ban đầu và đối chiếu với hiện tượng thực tế.
}
\end{ex}

% ===== Câu 119 =====
% ID: L11C2B10-30-TLN
% Mức: VD
\begin{ex}
% ID: L11C2B10-30-TLN
% Mức: VD
Từ kết quả của tình huống sau, nếu đại lượng cần tìm được nhân với hệ số 1.5 thì giá trị mới bằng bao nhiêu? Hai đỉnh thứ nhất và thứ chín cách nhau 16 ms. Tần số tính theo Hz bằng bao nhiêu?
\shortans{750}
\loigiai{
Xác định đúng dữ kiện và đơn vị, áp dụng hệ thức của dạng bài rồi tính được kết quả 750.
}
\end{ex}

% ===== Câu 120 =====
% ID: L11C2B10-30-TN
% Mức: VD
\dangbt{Thiết kế phương án thí nghiệm đo tần số của nguồn âm}
\begin{ex}
% ID: L11C2B10-30-TN
% Mức: VD
Trong chuyên đề “Thiết kế phương án thí nghiệm đo tần số của nguồn âm”, Phương pháp phù hợp nhất cho dạng “Thiết kế phương án thí nghiệm đo tần số của nguồn âm” là
\choice
{\True xác định dữ kiện, chọn đúng hệ thức hoặc định nghĩa, đổi đơn vị rồi kiểm tra kết quả}
{chỉ thay số mà không cần xét điều kiện áp dụng}
{bỏ qua đơn vị và chiều truyền sóng}
{luôn coi mọi điểm dao động cùng pha}
\loigiai{
Đáp án A. xác định dữ kiện, chọn đúng hệ thức hoặc định nghĩa, đổi đơn vị rồi kiểm tra kết quả. Đối chiếu định nghĩa, hệ thức hoặc dữ kiện của bài toán cho kết luận trên.
}
\end{ex}

% ===== Câu 121 =====
% ID: L11C2B10-31-TL
% Mức: TH
\dangbt{Đánh giá kết quả và sai số phép đo tần số âm}
\begin{ex}
% ID: L11C2B10-31-TL
% Mức: TH
Phân biệt sai số ngẫu nhiên và sai số hệ thống trong phép đo tần số âm, kèm một ví dụ cho mỗi loại.
\choiceTF

\loigiai{
Sai số ngẫu nhiên thay đổi giữa các lần đọc, ví dụ chọn vị trí đỉnh; sai số hệ thống lệch cùng chiều, ví dụ trục thời gian của thiết bị chưa hiệu chuẩn.
}
\end{ex}

% ===== Câu 122 =====
% ID: L11C2B10-31-TLN
% Mức: TH
\begin{ex}
% ID: L11C2B10-31-TLN
% Mức: TH
Đo được tần số $495\,\text{Hz}$, giá trị chuẩn $500\,\text{Hz}$. Sai số tương đối tính theo phần trăm bằng bao nhiêu?
\shortans{1}
\loigiai{
Xác định đúng dữ kiện và đơn vị, áp dụng hệ thức của dạng bài rồi tính được kết quả 1.
}
\end{ex}

% ===== Câu 123 =====
% ID: L11C2B10-31-TN
% Mức: NB
\dangbt{Xử lí số liệu nhiều lần đo và báo cáo kết quả tần số}
\begin{ex}
% ID: L11C2B10-31-TN
% Mức: NB
Trong chuyên đề “Xử lí số liệu nhiều lần đo và báo cáo kết quả tần số”, Nguồn sai số ngẫu nhiên thường gặp khi đọc đồ thị âm là
\choice
{ghi đúng số chu kì}
{\True xác định vị trí đỉnh không hoàn toàn giống nhau}
{dùng công thức $f=1/T$}
{đổi ms sang s}
\loigiai{
Đáp án B. xác định vị trí đỉnh không hoàn toàn giống nhau. Đối chiếu định nghĩa, hệ thức hoặc dữ kiện của bài toán cho kết luận trên.
}
\end{ex}

% ===== Câu 124 =====
% ID: L11C2B10-32-TL
% Mức: TH
\dangbt{Đánh giá kết quả và sai số phép đo tần số âm}
\begin{ex}
% ID: L11C2B10-32-TL
% Mức: TH
Một nhóm đo $T=(2,00\pm0,04)$ ms. Ước lượng tần số và độ không đảm bảo tuyệt đối của tần số.
\choiceTF

\loigiai{
$f=500$ Hz; sai số tương đối xấp xỉ 2%, nên $\Delta f\approx10$ Hz. Kết quả $f=(500\pm10)$ Hz.
}
\end{ex}

% ===== Câu 125 =====
% ID: L11C2B10-32-TLN
% Mức: TH
\begin{ex}
% ID: L11C2B10-32-TLN
% Mức: TH
Ba kết quả $498\,\text{Hz}$, $500\,\text{Hz}$, $502\,\text{Hz}$. Giá trị trung bình tính theo Hz bằng bao nhiêu?
\shortans{500}
\loigiai{
Xác định đúng dữ kiện và đơn vị, áp dụng hệ thức của dạng bài rồi tính được kết quả 500.
}
\end{ex}

% ===== Câu 126 =====
% ID: L11C2B10-32-TN
% Mức: NB
\dangbt{Xử lí số liệu nhiều lần đo và báo cáo kết quả tần số}
\begin{ex}
% ID: L11C2B10-32-TN
% Mức: NB
Trong chuyên đề “Xử lí số liệu nhiều lần đo và báo cáo kết quả tần số”, Biện pháp hiệu quả để giảm sai số tương đối của chu kì là
\choice
{tăng tiếng ồn nền}
{đổi tần số liên tục}
{\True đo nhiều chu kì rồi chia}
{chỉ đo nửa chu kì}
\loigiai{
Đáp án C. đo nhiều chu kì rồi chia. Đối chiếu định nghĩa, hệ thức hoặc dữ kiện của bài toán cho kết luận trên.
}
\end{ex}

% ===== Câu 127 =====
% ID: L11C2B10-33-TL
% Mức: VD
\dangbt{Đánh giá kết quả và sai số phép đo tần số âm}
\begin{ex}
% ID: L11C2B10-33-TL
% Mức: VD
Đề xuất ba cải tiến để tăng độ tin cậy của phép đo tần số âm.
\choiceTF

\loigiai{
Giảm tiếng ồn, đo nhiều chu kì và nhiều lần, chọn thang thời gian phù hợp/hiệu chuẩn thiết bị, giữ nguồn và microphone ổn định.
}
\end{ex}

% ===== Câu 128 =====
% ID: L11C2B10-33-TLN
% Mức: TH
\begin{ex}
% ID: L11C2B10-33-TLN
% Mức: TH
Chu kì đo được 2,0 ms với sai số 0,1 ms. Sai số tương đối của chu kì tính theo phần trăm bằng bao nhiêu?
\shortans{5}
\loigiai{
Xác định đúng dữ kiện và đơn vị, áp dụng hệ thức của dạng bài rồi tính được kết quả 5.
}
\end{ex}

% ===== Câu 129 =====
% ID: L11C2B10-33-TN
% Mức: TH
\dangbt{Xử lí số liệu nhiều lần đo và báo cáo kết quả tần số}
\begin{ex}
% ID: L11C2B10-33-TN
% Mức: TH
Trong chuyên đề “Xử lí số liệu nhiều lần đo và báo cáo kết quả tần số”, Nếu chu kì đo được lớn hơn giá trị thật thì tần số tính được sẽ
\choice
{lớn hơn giá trị thật}
{không đổi}
{luôn bằng 0}
{\True nhỏ hơn giá trị thật}
\loigiai{
Đáp án D. nhỏ hơn giá trị thật. Đối chiếu định nghĩa, hệ thức hoặc dữ kiện của bài toán cho kết luận trên.
}
\end{ex}

% ===== Câu 130 =====
% ID: L11C2B10-34-TL
% Mức: VD
\dangbt{Đánh giá kết quả và sai số phép đo tần số âm}
\begin{ex}
% ID: L11C2B10-34-TL
% Mức: VD
Phân tích các bước lập luận và chỉ ra điều kiện áp dụng trong bài toán: Phân biệt sai số ngẫu nhiên và sai số hệ thống trong phép đo tần số âm, kèm một ví dụ cho mỗi loại.
\choiceTF

\loigiai{
Sai số ngẫu nhiên thay đổi giữa các lần đọc, ví dụ chọn vị trí đỉnh; sai số hệ thống lệch cùng chiều, ví dụ trục thời gian của thiết bị chưa hiệu chuẩn. Cần nêu rõ điều kiện áp dụng, đổi đơn vị thống nhất và kiểm tra ý nghĩa vật lí của kết quả.
}
\end{ex}

% ===== Câu 131 =====
% ID: L11C2B10-34-TLN
% Mức: VD
\begin{ex}
% ID: L11C2B10-34-TLN
% Mức: VD
Từ kết quả của tình huống sau, nếu đại lượng cần tìm được nhân với hệ số 2 thì giá trị mới bằng bao nhiêu? Đo được tần số $495\,\text{Hz}$, giá trị chuẩn $500\,\text{Hz}$. Sai số tương đối tính theo phần trăm bằng bao nhiêu?
\shortans{2}
\loigiai{
Xác định đúng dữ kiện và đơn vị, áp dụng hệ thức của dạng bài rồi tính được kết quả 2.
}
\end{ex}

% ===== Câu 132 =====
% ID: L11C2B10-34-TN
% Mức: TH
\dangbt{Xử lí số liệu nhiều lần đo và báo cáo kết quả tần số}
\begin{ex}
% ID: L11C2B10-34-TN
% Mức: TH
Trong chuyên đề “Xử lí số liệu nhiều lần đo và báo cáo kết quả tần số”, Đo được tần số $495\,\text{Hz}$, giá trị chuẩn $500\,\text{Hz}$. Sai số tương đối tính theo phần trăm bằng bao nhiêu? Chọn giá trị đúng.
\choice
{\True 1}
{2}
{0.5}
{2}
\loigiai{
Đáp án A. 1. Đối chiếu định nghĩa, hệ thức hoặc dữ kiện của bài toán cho kết luận trên.
}
\end{ex}

% ===== Câu 133 =====
% ID: L11C2B10-35-TL
% Mức: VD
\dangbt{Đánh giá kết quả và sai số phép đo tần số âm}
\begin{ex}
% ID: L11C2B10-35-TL
% Mức: VD
Một học sinh giải bài sau nhưng bỏ qua điều kiện vật lí. Hãy sửa cách làm và trình bày lời giải đúng: Một nhóm đo $T=(2,00\pm0,04)$ ms. Ước lượng tần số và độ không đảm bảo tuyệt đối của tần số.
\choiceTF

\loigiai{
Cách giải đúng: $f=500$ Hz; sai số tương đối xấp xỉ 2%, nên $\Delta f\approx10$ Hz. Kết quả $f=(500\pm10)$ Hz. Sai lầm cần tránh là dùng hệ thức khi chưa xác định điều kiện biên, môi trường hoặc đơn vị.
}
\end{ex}

% ===== Câu 134 =====
% ID: L11C2B10-35-TLN
% Mức: VD
\begin{ex}
% ID: L11C2B10-35-TLN
% Mức: VD
Từ kết quả của tình huống sau, nếu đại lượng cần tìm được nhân với hệ số 0.5 thì giá trị mới bằng bao nhiêu? Ba kết quả $498\,\text{Hz}$, $500\,\text{Hz}$, $502\,\text{Hz}$. Giá trị trung bình tính theo Hz bằng bao nhiêu?
\shortans{250}
\loigiai{
Xác định đúng dữ kiện và đơn vị, áp dụng hệ thức của dạng bài rồi tính được kết quả 250.
}
\end{ex}

% ===== Câu 135 =====
% ID: L11C2B10-35-TN
% Mức: TH
\dangbt{Xử lí số liệu nhiều lần đo và báo cáo kết quả tần số}
\begin{ex}
% ID: L11C2B10-35-TN
% Mức: TH
Trong chuyên đề “Xử lí số liệu nhiều lần đo và báo cáo kết quả tần số”, Ba kết quả $498\,\text{Hz}$, $500\,\text{Hz}$, $502\,\text{Hz}$. Giá trị trung bình tính theo Hz bằng bao nhiêu? Chọn giá trị đúng.
\choice
{501}
{\True 500}
{1000}
{250}
\loigiai{
Đáp án B. 500. Đối chiếu định nghĩa, hệ thức hoặc dữ kiện của bài toán cho kết luận trên.
}
\end{ex}

% ===== Câu 136 =====
% ID: L11C2B10-36-TL
% Mức: VD
\dangbt{Đánh giá kết quả và sai số phép đo tần số âm}
\begin{ex}
% ID: L11C2B10-36-TL
% Mức: VD
Mở rộng tình huống sau thành một ví dụ thực tế và trình bày phương pháp giải: Đề xuất ba cải tiến để tăng độ tin cậy của phép đo tần số âm.
\choiceTF

\loigiai{
Giảm tiếng ồn, đo nhiều chu kì và nhiều lần, chọn thang thời gian phù hợp/hiệu chuẩn thiết bị, giữ nguồn và microphone ổn định. Có thể kiểm chứng bằng cách thay kết quả vào dữ kiện ban đầu và đối chiếu với hiện tượng thực tế.
}
\end{ex}

% ===== Câu 137 =====
% ID: L11C2B10-36-TLN
% Mức: VD
\begin{ex}
% ID: L11C2B10-36-TLN
% Mức: VD
Từ kết quả của tình huống sau, nếu đại lượng cần tìm được nhân với hệ số 1.5 thì giá trị mới bằng bao nhiêu? Chu kì đo được 2,0 ms với sai số 0,1 ms. Sai số tương đối của chu kì tính theo phần trăm bằng bao nhiêu?
\shortans{07/05/2026}
\loigiai{
Xác định đúng dữ kiện và đơn vị, áp dụng hệ thức của dạng bài rồi tính được kết quả 7.5.
}
\end{ex}

% ===== Câu 138 =====
% ID: L11C2B10-36-TN
% Mức: TH
\dangbt{Xử lí số liệu nhiều lần đo và báo cáo kết quả tần số}
\begin{ex}
% ID: L11C2B10-36-TN
% Mức: TH
Trong chuyên đề “Xử lí số liệu nhiều lần đo và báo cáo kết quả tần số”, Chu kì đo được 2,0 ms với sai số 0,1 ms. Sai số tương đối của chu kì tính theo phần trăm bằng bao nhiêu? Chọn giá trị đúng.
\choice
{02/05/2026}
{6}
{\True 5}
{10}
\loigiai{
Đáp án C. 5. Đối chiếu định nghĩa, hệ thức hoặc dữ kiện của bài toán cho kết luận trên.
}
\end{ex}

% ===== Câu 139 =====
% ID: L11C2B10-37-TN
% Mức: VD
\begin{ex}
% ID: L11C2B10-37-TN
% Mức: VD
Trong chuyên đề “Xử lí số liệu nhiều lần đo và báo cáo kết quả tần số”, Kết luận nào phù hợp nhất khi giải bài toán sau: Phân biệt sai số ngẫu nhiên và sai số hệ thống trong phép đo tần số âm, kèm một ví dụ cho mỗi loại.
\choice
{Không cần sử dụng định nghĩa hay dữ kiện đã cho.}
{Đại lượng cần tìm luôn bằng 0 trong mọi trường hợp.}
{Kết quả không phụ thuộc môi trường, tần số hoặc điều kiện biên.}
{\True Sai số ngẫu nhiên thay đổi giữa các lần đọc, ví dụ chọn vị trí đỉnh}
\loigiai{
Đáp án D. Sai số ngẫu nhiên thay đổi giữa các lần đọc, ví dụ chọn vị trí đỉnh. Đối chiếu định nghĩa, hệ thức hoặc dữ kiện của bài toán cho kết luận trên.
}
\end{ex}

% ===== Câu 140 =====
% ID: L11C2B10-38-TN
% Mức: VD
\begin{ex}
% ID: L11C2B10-38-TN
% Mức: VD
Trong chuyên đề “Xử lí số liệu nhiều lần đo và báo cáo kết quả tần số”, Kết luận nào phù hợp nhất khi giải bài toán sau: Một nhóm đo $T=(2,00\pm0,04)$ ms. Ước lượng tần số và độ không đảm bảo tuyệt đối của tần số.
\choice
{\True $f=500$ Hz}
{Không cần sử dụng định nghĩa hay dữ kiện đã cho.}
{Đại lượng cần tìm luôn bằng 0 trong mọi trường hợp.}
{Kết quả không phụ thuộc môi trường, tần số hoặc điều kiện biên.}
\loigiai{
Đáp án A. $f=500$ Hz. Đối chiếu định nghĩa, hệ thức hoặc dữ kiện của bài toán cho kết luận trên.
}
\end{ex}

% ===== Câu 141 =====
% ID: L11C2B10-39-TN
% Mức: VD
\begin{ex}
% ID: L11C2B10-39-TN
% Mức: VD
Trong chuyên đề “Xử lí số liệu nhiều lần đo và báo cáo kết quả tần số”, Kết luận nào phù hợp nhất khi giải bài toán sau: Đề xuất ba cải tiến để tăng độ tin cậy của phép đo tần số âm.
\choice
{Kết quả không phụ thuộc môi trường, tần số hoặc điều kiện biên.}
{\True Giảm tiếng ồn, đo nhiều chu kì và nhiều lần, chọn thang thời gian phù hợp/hiệu chuẩn thiết bị, giữ nguồn và microphone ổn định.}
{Không cần sử dụng định nghĩa hay dữ kiện đã cho.}
{Đại lượng cần tìm luôn bằng 0 trong mọi trường hợp.}
\loigiai{
Đáp án B. Giảm tiếng ồn, đo nhiều chu kì và nhiều lần, chọn thang thời gian phù hợp/hiệu chuẩn thiết bị, giữ nguồn và microphone ổn định.. Đối chiếu định nghĩa, hệ thức hoặc dữ kiện của bài toán cho kết luận trên.
}
\end{ex}

% ===== Câu 142 =====
% ID: L11C2B10-40-TN
% Mức: VD
\begin{ex}
% ID: L11C2B10-40-TN
% Mức: VD
Trong chuyên đề “Xử lí số liệu nhiều lần đo và báo cáo kết quả tần số”, Phương pháp phù hợp nhất cho dạng “Xử lí số liệu nhiều lần đo và báo cáo kết quả tần số” là
\choice
{bỏ qua đơn vị và chiều truyền sóng}
{luôn coi mọi điểm dao động cùng pha}
{\True xác định dữ kiện, chọn đúng hệ thức hoặc định nghĩa, đổi đơn vị rồi kiểm tra kết quả}
{chỉ thay số mà không cần xét điều kiện áp dụng}
\loigiai{
Đáp án C. xác định dữ kiện, chọn đúng hệ thức hoặc định nghĩa, đổi đơn vị rồi kiểm tra kết quả. Đối chiếu định nghĩa, hệ thức hoặc dữ kiện của bài toán cho kết luận trên.
}
\end{ex}

% ===== Câu 143 =====
% ID: L11C2B10-41-TN
% Mức: NB
\dangbt{Xử lí đồ thị dao động âm để xác định chu kì và tần số}
\begin{ex}
% ID: L11C2B10-41-TN
% Mức: NB
Trên đồ thị, 10 chu kì kéo dài 20 ms. Tần số âm bằng
\choice
{$2000\,\text{Hz}$}
{\True $500\,\text{Hz}$}
{$50\,\text{Hz}$}
{$200\,\text{Hz}$}
\loigiai{
Đáp án B. $500\,\text{Hz}$. Đối chiếu định nghĩa, hệ thức hoặc dữ kiện của bài toán cho kết luận trên.
}
\end{ex}

% ===== Câu 144 =====
% ID: L11C2B10-42-TN
% Mức: NB
\begin{ex}
% ID: L11C2B10-42-TN
% Mức: NB
Khoảng thời gian giữa hai đỉnh liên tiếp là 2,5 ms. Tần số bằng
\choice
{$40\,\text{Hz}$}
{4 kHz}
{\True $400\,\text{Hz}$}
{$250\,\text{Hz}$}
\loigiai{
Đáp án C. $400\,\text{Hz}$. Đối chiếu định nghĩa, hệ thức hoặc dữ kiện của bài toán cho kết luận trên.
}
\end{ex}

% ===== Câu 145 =====
% ID: L11C2B10-43-TN
% Mức: TH
\begin{ex}
% ID: L11C2B10-43-TN
% Mức: TH
Để giảm sai số đọc chu kì từ đồ thị, nên
\choice
{chỉ chọn một đoạn rất ngắn}
{đọc biên độ thay cho thời gian}
{thay đổi tần số khi đang đo}
{\True đo thời gian của nhiều chu kì}
\loigiai{
Đáp án D. đo thời gian của nhiều chu kì. Đối chiếu định nghĩa, hệ thức hoặc dữ kiện của bài toán cho kết luận trên.
}
\end{ex}

% ===== Câu 146 =====
% ID: L11C2B10-44-TN
% Mức: TH
\begin{ex}
% ID: L11C2B10-44-TN
% Mức: TH
12 chu kì kéo dài 24 ms. Tần số tính theo Hz bằng bao nhiêu? Chọn giá trị đúng.
\choice
{\True 500}
{1000}
{250}
{501}
\loigiai{
Đáp án A. 500. Đối chiếu định nghĩa, hệ thức hoặc dữ kiện của bài toán cho kết luận trên.
}
\end{ex}

% ===== Câu 147 =====
% ID: L11C2B10-45-TN
% Mức: TH
\begin{ex}
% ID: L11C2B10-45-TN
% Mức: TH
Khoảng thời gian 4,5 ms chứa 3 chu kì. Tần số tính theo Hz, làm tròn đến hàng đơn vị, bằng bao nhiêu? Chọn giá trị đúng.
\choice
{668}
{\True 667}
{1334}
{333.5}
\loigiai{
Đáp án B. 667. Đối chiếu định nghĩa, hệ thức hoặc dữ kiện của bài toán cho kết luận trên.
}
\end{ex}

% ===== Câu 148 =====
% ID: L11C2B10-46-TN
% Mức: TH
\begin{ex}
% ID: L11C2B10-46-TN
% Mức: TH
Hai đỉnh thứ nhất và thứ chín cách nhau 16 ms. Tần số tính theo Hz bằng bao nhiêu? Chọn giá trị đúng.
\choice
{250}
{501}
{\True 500}
{1000}
\loigiai{
Đáp án C. 500. Đối chiếu định nghĩa, hệ thức hoặc dữ kiện của bài toán cho kết luận trên.
}
\end{ex}

% ===== Câu 149 =====
% ID: L11C2B10-47-TN
% Mức: VD
\begin{ex}
% ID: L11C2B10-47-TN
% Mức: VD
Kết luận nào phù hợp nhất khi giải bài toán sau: Đồ thị âm cho thấy từ đỉnh thứ 2 đến đỉnh thứ 12 là 25 ms. Tính chu kì và tần số.
\choice
{Không cần sử dụng định nghĩa hay dữ kiện đã cho.}
{Đại lượng cần tìm luôn bằng 0 trong mọi trường hợp.}
{Kết quả không phụ thuộc môi trường, tần số hoặc điều kiện biên.}
{\True Giữa 11 đỉnh có 10 chu kì: $T=25/10=2,5$ ms}
\loigiai{
Đáp án D. Giữa 11 đỉnh có 10 chu kì: $T=25/10=2,5$ ms. Đối chiếu định nghĩa, hệ thức hoặc dữ kiện của bài toán cho kết luận trên.
}
\end{ex}

% ===== Câu 150 =====
% ID: L11C2B10-48-TN
% Mức: VD
\begin{ex}
% ID: L11C2B10-48-TN
% Mức: VD
Kết luận nào phù hợp nhất khi giải bài toán sau: Ba lần đo thời gian của 10 chu kì là 19,8 ms; 20,0 ms; 20,2 ms. Tính tần số từ giá trị trung bình.
\choice
{\True Thời gian trung bình của 10 chu kì là 20,0 ms}
{Không cần sử dụng định nghĩa hay dữ kiện đã cho.}
{Đại lượng cần tìm luôn bằng 0 trong mọi trường hợp.}
{Kết quả không phụ thuộc môi trường, tần số hoặc điều kiện biên.}
\loigiai{
Đáp án A. Thời gian trung bình của 10 chu kì là 20,0 ms. Đối chiếu định nghĩa, hệ thức hoặc dữ kiện của bài toán cho kết luận trên.
}
\end{ex}

% ===== Câu 151 =====
% ID: L11C2B10-49-TN
% Mức: VD
\begin{ex}
% ID: L11C2B10-49-TN
% Mức: VD
Kết luận nào phù hợp nhất khi giải bài toán sau: Mô tả cách nhận biết và loại bỏ một đoạn tín hiệu bị nhiễu trước khi tính tần số.
\choice
{Kết quả không phụ thuộc môi trường, tần số hoặc điều kiện biên.}
{\True Chọn đoạn tuần hoàn ổn định, biên dạng lặp rõ}
{Không cần sử dụng định nghĩa hay dữ kiện đã cho.}
{Đại lượng cần tìm luôn bằng 0 trong mọi trường hợp.}
\loigiai{
Đáp án B. Chọn đoạn tuần hoàn ổn định, biên dạng lặp rõ. Đối chiếu định nghĩa, hệ thức hoặc dữ kiện của bài toán cho kết luận trên.
}
\end{ex}

% ===== Câu 152 =====
% ID: L11C2B10-50-TN
% Mức: VD
\begin{ex}
% ID: L11C2B10-50-TN
% Mức: VD
Phương pháp phù hợp nhất cho dạng “Xử lí đồ thị dao động âm để xác định chu kì và tần số” là
\choice
{bỏ qua đơn vị và chiều truyền sóng}
{luôn coi mọi điểm dao động cùng pha}
{\True xác định dữ kiện, chọn đúng hệ thức hoặc định nghĩa, đổi đơn vị rồi kiểm tra kết quả}
{chỉ thay số mà không cần xét điều kiện áp dụng}
\loigiai{
Đáp án C. xác định dữ kiện, chọn đúng hệ thức hoặc định nghĩa, đổi đơn vị rồi kiểm tra kết quả. Đối chiếu định nghĩa, hệ thức hoặc dữ kiện của bài toán cho kết luận trên.
}
\end{ex}

% ===== Câu 153 =====
% ID: L11C2B10-51-TN
% Mức: NB
\dangbt{Đánh giá kết quả và sai số phép đo tần số âm}
\begin{ex}
% ID: L11C2B10-51-TN
% Mức: NB
Nguồn sai số ngẫu nhiên thường gặp khi đọc đồ thị âm là
\choice
{đổi ms sang s}
{ghi đúng số chu kì}
{\True xác định vị trí đỉnh không hoàn toàn giống nhau}
{dùng công thức $f=1/T$}
\loigiai{
Đáp án C. xác định vị trí đỉnh không hoàn toàn giống nhau. Đối chiếu định nghĩa, hệ thức hoặc dữ kiện của bài toán cho kết luận trên.
}
\end{ex}

% ===== Câu 154 =====
% ID: L11C2B10-52-TN
% Mức: NB
\begin{ex}
% ID: L11C2B10-52-TN
% Mức: NB
Biện pháp hiệu quả để giảm sai số tương đối của chu kì là
\choice
{chỉ đo nửa chu kì}
{tăng tiếng ồn nền}
{đổi tần số liên tục}
{\True đo nhiều chu kì rồi chia}
\loigiai{
Đáp án D. đo nhiều chu kì rồi chia. Đối chiếu định nghĩa, hệ thức hoặc dữ kiện của bài toán cho kết luận trên.
}
\end{ex}

% ===== Câu 155 =====
% ID: L11C2B10-53-TN
% Mức: TH
\begin{ex}
% ID: L11C2B10-53-TN
% Mức: TH
Nếu chu kì đo được lớn hơn giá trị thật thì tần số tính được sẽ
\choice
{\True nhỏ hơn giá trị thật}
{lớn hơn giá trị thật}
{không đổi}
{luôn bằng 0}
\loigiai{
Đáp án A. nhỏ hơn giá trị thật. Đối chiếu định nghĩa, hệ thức hoặc dữ kiện của bài toán cho kết luận trên.
}
\end{ex}

% ===== Câu 156 =====
% ID: L11C2B10-54-TN
% Mức: TH
\begin{ex}
% ID: L11C2B10-54-TN
% Mức: TH
Đo được tần số $495\,\text{Hz}$, giá trị chuẩn $500\,\text{Hz}$. Sai số tương đối tính theo phần trăm bằng bao nhiêu? Chọn giá trị đúng.
\choice
{2}
{\True 1}
{2}
{0.5}
\loigiai{
Đáp án B. 1. Đối chiếu định nghĩa, hệ thức hoặc dữ kiện của bài toán cho kết luận trên.
}
\end{ex}

% ===== Câu 157 =====
% ID: L11C2B10-55-TN
% Mức: TH
\begin{ex}
% ID: L11C2B10-55-TN
% Mức: TH
Ba kết quả $498\,\text{Hz}$, $500\,\text{Hz}$, $502\,\text{Hz}$. Giá trị trung bình tính theo Hz bằng bao nhiêu? Chọn giá trị đúng.
\choice
{250}
{501}
{\True 500}
{1000}
\loigiai{
Đáp án C. 500. Đối chiếu định nghĩa, hệ thức hoặc dữ kiện của bài toán cho kết luận trên.
}
\end{ex}

% ===== Câu 158 =====
% ID: L11C2B10-56-TN
% Mức: TH
\begin{ex}
% ID: L11C2B10-56-TN
% Mức: TH
Chu kì đo được 2,0 ms với sai số 0,1 ms. Sai số tương đối của chu kì tính theo phần trăm bằng bao nhiêu? Chọn giá trị đúng.
\choice
{10}
{02/05/2026}
{6}
{\True 5}
\loigiai{
Đáp án D. 5. Đối chiếu định nghĩa, hệ thức hoặc dữ kiện của bài toán cho kết luận trên.
}
\end{ex}

% ===== Câu 159 =====
% ID: L11C2B10-57-TN
% Mức: VD
\begin{ex}
% ID: L11C2B10-57-TN
% Mức: VD
Kết luận nào phù hợp nhất khi giải bài toán sau: Phân biệt sai số ngẫu nhiên và sai số hệ thống trong phép đo tần số âm, kèm một ví dụ cho mỗi loại.
\choice
{\True Sai số ngẫu nhiên thay đổi giữa các lần đọc, ví dụ chọn vị trí đỉnh}
{Không cần sử dụng định nghĩa hay dữ kiện đã cho.}
{Đại lượng cần tìm luôn bằng 0 trong mọi trường hợp.}
{Kết quả không phụ thuộc môi trường, tần số hoặc điều kiện biên.}
\loigiai{
Đáp án A. Sai số ngẫu nhiên thay đổi giữa các lần đọc, ví dụ chọn vị trí đỉnh. Đối chiếu định nghĩa, hệ thức hoặc dữ kiện của bài toán cho kết luận trên.
}
\end{ex}

% ===== Câu 160 =====
% ID: L11C2B10-58-TN
% Mức: VD
\begin{ex}
% ID: L11C2B10-58-TN
% Mức: VD
Kết luận nào phù hợp nhất khi giải bài toán sau: Một nhóm đo $T=(2,00\pm0,04)$ ms. Ước lượng tần số và độ không đảm bảo tuyệt đối của tần số.
\choice
{Kết quả không phụ thuộc môi trường, tần số hoặc điều kiện biên.}
{\True $f=500$ Hz}
{Không cần sử dụng định nghĩa hay dữ kiện đã cho.}
{Đại lượng cần tìm luôn bằng 0 trong mọi trường hợp.}
\loigiai{
Đáp án B. $f=500$ Hz. Đối chiếu định nghĩa, hệ thức hoặc dữ kiện của bài toán cho kết luận trên.
}
\end{ex}

% ===== Câu 161 =====
% ID: L11C2B10-59-TN
% Mức: VD
\begin{ex}
% ID: L11C2B10-59-TN
% Mức: VD
Kết luận nào phù hợp nhất khi giải bài toán sau: Đề xuất ba cải tiến để tăng độ tin cậy của phép đo tần số âm.
\choice
{Đại lượng cần tìm luôn bằng 0 trong mọi trường hợp.}
{Kết quả không phụ thuộc môi trường, tần số hoặc điều kiện biên.}
{\True Giảm tiếng ồn, đo nhiều chu kì và nhiều lần, chọn thang thời gian phù hợp/hiệu chuẩn thiết bị, giữ nguồn và microphone ổn định.}
{Không cần sử dụng định nghĩa hay dữ kiện đã cho.}
\loigiai{
Đáp án C. Giảm tiếng ồn, đo nhiều chu kì và nhiều lần, chọn thang thời gian phù hợp/hiệu chuẩn thiết bị, giữ nguồn và microphone ổn định.. Đối chiếu định nghĩa, hệ thức hoặc dữ kiện của bài toán cho kết luận trên.
}
\end{ex}

% ===== Câu 162 =====
% ID: L11C2B10-60-TN
% Mức: VD
\begin{ex}
% ID: L11C2B10-60-TN
% Mức: VD
Phương pháp phù hợp nhất cho dạng “Đánh giá kết quả và sai số phép đo tần số âm” là
\choice
{chỉ thay số mà không cần xét điều kiện áp dụng}
{bỏ qua đơn vị và chiều truyền sóng}
{luôn coi mọi điểm dao động cùng pha}
{\True xác định dữ kiện, chọn đúng hệ thức hoặc định nghĩa, đổi đơn vị rồi kiểm tra kết quả}
\loigiai{
Đáp án D. xác định dữ kiện, chọn đúng hệ thức hoặc định nghĩa, đổi đơn vị rồi kiểm tra kết quả. Đối chiếu định nghĩa, hệ thức hoặc dữ kiện của bài toán cho kết luận trên.
}
\end{ex}
